\documentclass[11pt]{amsart}
\usepackage[margin=1.05in]{geometry}
\usepackage{amsmath,amssymb,amsthm,mathtools}
\usepackage{booktabs}
\usepackage{longtable}
\usepackage{array}
\usepackage{xcolor}
\usepackage[hidelinks]{hyperref}
\usepackage{url}
\usepackage{xurl}
\usepackage[T1]{fontenc}
\sloppy
\allowdisplaybreaks

\newtheorem{theorem}{Theorem}[section]
\newtheorem{proposition}[theorem]{Proposition}
\newtheorem{lemma}[theorem]{Lemma}
\newtheorem{corollary}[theorem]{Corollary}
\newtheorem{conjecture}[theorem]{Conjecture}
\newtheorem*{theoremA}{Theorem A}
\newtheorem*{theoremB}{Theorem B}
\newtheorem*{theoremC}{Theorem C}
\newtheorem*{theoremD}{Theorem D}
\newtheorem*{theoremE}{Theorem E}
\newtheorem*{theoremF}{Theorem F}
\newtheorem*{theoremG}{Theorem G}
\theoremstyle{definition}
\newtheorem{definition}[theorem]{Definition}
\newtheorem{problem}[theorem]{Problem}
\theoremstyle{remark}
\newtheorem{remark}[theorem]{Remark}
\newtheorem{observation}[theorem]{Observation}
\newtheorem{example}[theorem]{Example}

\newcommand{\lr}{\ell}
\newcommand{\pp}{p}
\newcommand{\CA}{\textup{[computer-assisted]}}
\newcommand{\GEN}{\mathrm{GEN}}
\newcommand{\RAD}{\mathrm{RAD}}
\newcommand{\ROOT}{\mathrm{ROOT}}

\title[On Graffiti/WOWII Conjecture 133]{On Graffiti/WOWII Conjecture 133: small neighbourhood independence, a reduction to bounded induced paths, and computer-assisted bounds}

\author[Neil Xu]{Neil Xu\\[2pt]{\normalfont\small\texttt{\lowercase{neard.ws@gmail.com}}}}

\date{October 2026}

\begin{document}

\begin{abstract}
Conjecture~133 of \emph{Written on the Wall II}, generated by the program Graffiti.pc, states that every connected graph $G$ on at least two vertices with no $4$-cycle subgraph satisfies $\pp(G)\ge r(G)+\lr(G)$. Here $\pp(G)$ is the largest order of an induced path, $r(G)$ is the radius, and $\lr(G)=S(G)/n$ is the average over the vertices $v$ of $\alpha(G[N(v)])$. We prove the conjecture for every graph with $\lr(G)\le5$. Using a theorem of Erd\H{o}s, Saks and S\'os, the conjecture reduces to two families of statements: radius-free mass bounds $\GEN_K$ (no induced path on $2K+1$ vertices implies $S(G)\le Kn$), and radius-extremal bounds $\RAD(r)$ (radius $r$ and $\pp=2r-1$ imply $\lr\le r-1$). We prove $\GEN_3$ by hand together with three finite local certificates, and we have formalised this proof completely in Lean~4. We prove $\GEN_4$ with a finite extension certificate. For the radius-extremal side we introduce rooted mass bounds $\ROOT(q)$: if every induced path starting at a root $s$ has at most $q$ vertices, then $S\le(q-1)(n-1)-(q-3)d(s)-m$, where $m$ is the number of edges in $N(s)$. We prove $\ROOT(q)$ for $3\le q\le8$ and $\RAD(r)$ for $4\le r\le10$, the latter in the sharp form $S\le(r-1)n-2r(r-3)$ for $5\le r\le 10$, with equality only for $C_{2r}$. Both families are proved for all indices in two restricted classes of strong $C_4$-free graphs: bipartite graphs, and graphs whose blocks are bipartite or cycles (the radius bounds in the second class use a structure theorem for radius-extremal graphs). For $5<\lr\le6$ every counterexample has $(r,\pp)\in\{(4,9),(5,10)\}$; we close the bipartite case and several further subcases of these two cells, but not the cells themselves. We also give counterexamples to several natural strengthenings: a star bound, an endpoint inequality, a dominating-core budget (refuted by the Petersen graph), and one specific scalar relaxation (refuted by data of the Heawood graph). Most of the new results are computer-assisted. Each was produced in one session of an AI model and then checked by a separate session of the same model acting as referee, which recomputed the finite certificates with its own code; this is not external peer review. Every statement is labelled with its status. In a continuous computer-assisted proving window ending 24~October~2026 (Asia/Shanghai), rounds R35--R40 produced further accepted partial mass theorems toward $\mathrm{EB}_6$ and triangle-free $\GEN_6$ (high-vertex bounds, endpoint-cap deficits, Robertson puncture classification, and conditional residual equivalences); universal $\mathrm{EB}_6$, unrestricted triangle-free $\GEN_6$, and REAL-133 itself remain open.
\end{abstract}

\maketitle

\section{Introduction}

All graphs are finite and simple. For a graph $G$ with $n\ge2$ vertices let $\pp(G)$ be the largest number of vertices of an induced path, $r(G)$ the radius (measured in edges), and
\[
\lr(G)=\frac1n\sum_{v\in V(G)}\alpha\bigl(G[N(v)]\bigr),
\]
where $\alpha$ is the independence number and $N(v)$ the open neighbourhood. Write $c_{C_4}(G)=1$ if $G$ contains no $4$-cycle as a subgraph (chorded or not) and $c_{C_4}(G)=0$ otherwise.

\begin{conjecture}[Graffiti/WOWII Conjecture 133]\label{conj:133}
Every connected graph $G$ on at least two vertices satisfies
\[
r(G)+\lr(G)^{\,c_{C_4}(G)}\le \pp(G).
\]
\end{conjecture}

The conjecture is number~133 in E.~DeLaVina's list \emph{Written on the Wall II} of conjectures of the program Graffiti.pc~\cite{wowii}. We follow the formal statement in the Formal Conjectures repository~\cite{fc133}, file \texttt{FormalConjectures/WrittenOnTheWallII/GraphConjecture133.lean}, which reads \texttt{(rad : $\mathbb{R}$) + l G \^{} cC4 $\le$ (path G : $\mathbb{R}$)}. Earlier versions of that file had $\lfloor \lr(G)\rfloor$ in place of $\lr(G)$. Pull request \#5727 of the repository~\cite{fcpr}, ``drop the floor on \texttt{l G}'' (merged on 13~September 2026), removed the floor on the grounds that the brackets in the WOWII source denote grouping. We work with the real-valued statement throughout and call it REAL-133.

If $G$ contains a $4$-cycle, the conjecture reads $\pp(G)\ge r(G)+1$, and this always holds: a shortest path between two vertices at distance $\operatorname{diam}(G)\ge r(G)$ is induced. So only graphs without a $4$-cycle subgraph matter. We call them \emph{strong $C_4$-free}; the word ``strong'' emphasises that chorded $4$-cycles are excluded too. Since $\pp$ and $r$ are integers, for a strong $C_4$-free $G$ the conjecture is equivalent to
\begin{equation}\label{eq:target}
\pp(G)-r(G)\ \ge\ \lceil \lr(G)\rceil .
\end{equation}

H.~Chen~\cite{chen} has studied mainly the earlier floor version $\pp\ge r+\lfloor\lr\rfloor$. That preprint describes itself as AI-assisted and not peer-reviewed, and we have not checked its proofs. For the floor version it claims proofs when $\lfloor\lr\rfloor\le2$, when $\lfloor\lr\rfloor=3$ and $r=2$, and for minimum-order counterexamples with $\lfloor\lr\rfloor=3$ and $r=3$. For the real-valued form it records an exhaustive check on strong $C_4$-free graphs with at most $13$ vertices, and its Theorem~6.7 states that $\operatorname{diam}(G)=2$ and $\lr>3$ imply $\pp\ge6$; this overlaps with the case $r=2$, $3<\lr\le4$ of our Theorem~A for graphs of diameter $2$. Its Lemma~7.18 contains the trace classification in the first part of our Lemma~\ref{lem:c6}. Our proofs are independent of~\cite{chen}.

\subsection*{Status labels}
Every statement carries one of the following labels.
\begin{itemize}
\item \textbf{Proved}: a complete written proof in this paper, possibly using the published theorems quoted in Section~\ref{sec:prelim}.
\item \textbf{Proved (written proof in the supplement)}: a complete written proof without computer assistance that is too long to reproduce here; the exact file is named in the statement and in Appendix~\ref{app:index}.
\item \textbf{Finite classification}: a written proof that also uses an exhaustive classification of small base configurations (at most $13$ vertices; single-vertex attachment tests on them have at most $14$ vertices), re-run in seconds by the included stand-alone checker.
\item \textbf{Computer-assisted}: the proof consists of written reductions together with finite certificates (extension trees, trace tables, exclusion lists) whose checking requires a computer. For these results the paper gives the statement, the structure of the argument and the size of the certificates; the complete arguments and certificates are in the supplementary material (Section~\ref{sec:data}). Every such result passed the verification protocol of Section~\ref{sec:method}; Appendix~\ref{app:index} lists, for each, the written proof, the certificate files and the referee report.
\item \textbf{Formalised}: the statement and its proof are checked by the Lean~4 kernel.
\item \textbf{Conjecture} or \textbf{open}: not proved. Finite evidence, if any, is stated with its exact domain.
\end{itemize}

\subsection*{The reduction}
Write $S(G)=\sum_v\alpha(G[N(v)])=n\lr(G)$. For integers $K\ge3$ and $r\ge4$ consider:
\begin{itemize}
\item[$\GEN_K$:] every connected strong $C_4$-free graph $G$ with $\pp(G)\le2K$ satisfies $S(G)\le Kn$;
\item[$\RAD(r)$:] every connected strong $C_4$-free graph with $r(G)=r$ and $\pp(G)=2r-1$ satisfies $\lr(G)\le r-1$.
\end{itemize}

\begin{theoremB}[proved; Theorem~\ref{thm:B}]
Let $G$ be a counterexample to REAL-133. Then $\lr(G)>3$, and with $K=\lceil\lr(G)\rceil-1\ge3$, either $r(G)\le K$ and $\pp(G)\le2K$ \textup{(}so $G$ violates $\GEN_K$\textup{)}, or $r(G)=K+1$ and $\pp(G)=2r(G)-1$ \textup{(}so $G$ violates $\RAD(K+1)$\textup{)}.
\end{theoremB}

\subsection*{Main results}

\begin{theoremA}[proved; Theorem~\ref{thm:A}]
REAL-133 holds for every connected strong $C_4$-free graph $G$ with $\lr(G)\le3$, and for every one with $r(G)\le2$ and $\lr(G)\le 4$.
\end{theoremA}

\begin{theoremC}[$\GEN_3$ proved with three finite certificates, and formalised; $\GEN_4$ computer-assisted]
Every connected strong $C_4$-free graph with no induced $P_7$ satisfies $S(G)\le3n$ \textup{(}Theorem~\ref{thm:gen3}; formalised in Lean~4, Section~\ref{sec:lean}\textup{)}. Every connected strong $C_4$-free graph with no induced $P_9$ satisfies $S(G)\le4n$ \textup{(}Theorem~\ref{thm:gen4}\textup{)}.
\end{theoremC}

For a vertex $s$ of a connected graph $F$ let $\pp_s(F)$ be the largest order of an induced path starting at $s$, $D=d_F(s)$, and $m=e(F[N(s)])$. Consider
\begin{itemize}
\item[$\ROOT(q)$:] every connected strong $C_4$-free $F$ with a vertex $s$ such that $\pp_s(F)\le q$ satisfies
\[
S(F)\le(q-1)(|F|-1)-(q-3)D-m .
\]
\end{itemize}

\begin{theoremD}[computer-assisted; Sections~\ref{sec:rad}--\ref{sec:root}]
\begin{enumerate}
\item $\ROOT(q)$ holds for $3\le q\le 8$, and its coefficients $q-3$ and $1$ are optimal for the leading coefficient $q-1$.
\item If $G$ is connected and strong $C_4$-free with $r(G)=4$ and $\pp(G)=7$, then $S(G)\le2n$ \textup{(}finite classification\textup{)}; with $r(G)=5$ and $\pp(G)=9$, $S(G)=2n-T(G)$.
\item For $5\le r\le10$, every connected strong $C_4$-free $G$ with $r(G)=r$ and $\pp(G)=2r-1$ satisfies
\[
S(G)\le(r-1)n-2r(r-3),
\]
with equality only for the cycle $C_{2r}$. In particular $\RAD(r)$ holds for $4\le r\le10$.
\end{enumerate}
\end{theoremD}

\begin{theoremE}[consequence]
REAL-133 holds for every connected graph with $\lr(G)\le5$. A counterexample with $5<\lr(G)\le6$ has $(r(G),\pp(G))\in\{(4,9),(5,10)\}$. For every $K\le 9$, alternative~\textup{(b)} of Theorem~B cannot occur.
\end{theoremE}

\begin{theoremF}[restricted classes; proved, partly computer-assisted]
$\ROOT(q)$ for every $q$, and the sharp bound of Theorem~D(3) for every $r\ge5$, hold for every connected strong $C_4$-free graph all of whose blocks are bipartite or cycles \textup{(}Theorem~\ref{thm:blocks}; the radius statement uses $\mathrm{Struct}(r)$\textup{)}. Every bipartite strong $C_4$-free graph with no induced $P_{11}$ satisfies $S\le5n-8-|a-b|$, where $a,b$ are the sizes of the colour classes \textup{(}Theorem~\ref{thm:bip}; computer-assisted\textup{)}; so $\GEN_5$ holds for bipartite graphs.
\end{theoremF}

Section~\ref{sec:neg} collects negative results that rule out several natural routes: $\GEN_K$ cannot be proved from singleton deletion losses (star bounds fail for $K=5$); a natural endpoint inequality fails at $q=9$; a dominating-core budget suggested by the computations is refuted by the Petersen graph; and the data of the Heawood graph satisfy one specific system of scalar constraints (degree, triangle and distance identities together with the two generic path lower bounds) while violating the conjectured inequality, so that this particular system cannot imply REAL-133. Grafts of projective-plane incidence graphs show that $\lr$ is unbounded among radius-extremal graphs ($\pp=2r-1$), so $\RAD(r)$ must depend on~$r$. Section~\ref{sec:open} states what is open, including two conjectures (TD and $\mathrm{CB}_{\mathrm{high}}$), the first of which would settle the radius-extremal side of the reduction and the second the whole conjecture. Partial results on the remaining band $5<\lr\le6$ are collected in Appendix~\ref{app:cells}.

\begin{theoremG}[computer-assisted; continuous window R35--R40; Appendix~\ref{app:window}]
In the triangle-free strong $C_4$-free class with longest induced path order $p\le12$, the following were accepted by separate Astra referee sessions as \textsc{correct} or \textsc{correct with fixes} (packaging fixes only). Write $S$ for the surplus and call a set $X$ \emph{$6$-cheap} if the deletion loss satisfies $D(X)\le6|X|$.
\begin{enumerate}
\item \textbf{Mass and deficits.} Sharp deficit $S\le6n-28$ for $p\le11$, $n\ge8$; deficit $S\le6n-36$ under $P_{12}$-endpoint degree cap four and $n\ge14$; endpoint-reducing deletions of weight at most $18$ leaving at least $14$ vertices imply $S\le6n$. High-vertex mass: at most $31$ (resp.\ $32$) vertices of degree $\ge5$ imply $S\le6n$; under $P_{12}$-endpoint cap five the bound extends to $35$ (resp.\ $36$) high vertices. Cap-four endpoint systems and listed degree-$5/6/7$ contact families closed in R35--R38 are retained as accepted background.
\item \textbf{Structural tools.} Rooted $P_7$ and $2$-connectivity lemmas; pair/penalty identities for degree-four forests; explicit critical examples $C_{106}$ and $C_{112}$ (minimum high-vertex count $29$ for a nonempty graph with no $6$-cheap set). Every nonempty triangle-free strong $C_4$-free graph of minimum degree $\ge4$ contains an induced $P_{11}$ (equivalently $p\le10$ is $3$-degenerate).
\item \textbf{Robertson punctures.} A connected minimum-degree-$4$ graph containing a one-vertex puncture of the Robertson graph is isomorphic to $R$, $X_{21}$, $X_{22}$ or $X_{23}$; whenever such a puncture is induced there is a nonempty $6$-cheap set of order $\le7$ (sharp on $X_{22}$); every positive graph at this path cap reduces to a positive induced subgraph free of all three punctures.
\item \textbf{Conditional assembly.} Relative to the accepted background, cheap deletion on the stated residual is equivalent to universal $\mathrm{EB}_6$ and to unrestricted triangle-free $\GEN_6$; every proved $6$-cheap set is $K$-cheap for real $K\ge6$ at the unchanged path cap $12$. No emptiness of the avoiding class A38, and no unrestricted $\GEN_6$, is claimed.
\end{enumerate}
\end{theoremG}


\begin{table}[ht]
\centering\small
\begin{tabular}{@{}p{6.4cm}p{4.6cm}p{2.6cm}@{}}
\toprule
Statement & Status & Where\\
\midrule
REAL-133 for $\lr\le3$; for $r\le2$, $\lr\le4$ & proved & Thm.~\ref{thm:A}\\
Reduction to $\GEN_K$ and $\RAD(r)$ & proved & Thm.~\ref{thm:B}\\
$\GEN_3$ & proved + 3 finite certificates; formalised & Thm.~\ref{thm:gen3}, \S\ref{sec:lean}\\
$\GEN_4$ & computer-assisted & Thm.~\ref{thm:gen4}\\
$\RAD(4)$: $r=4,\pp=7\Rightarrow S\le2n$ & finite classification & Thm.~\ref{thm:r4}\\
$\RAD(5)$: $S=2n-T$ & computer-assisted & Thm.~\ref{thm:rad5}\\
$\RAD(6)$, $\RAD(7)$ & computer-assisted & Thms.~\ref{thm:rad6}, \ref{thm:rad7}\\
$\ROOT(3)$; $\ROOT(4)$, $\ROOT(5)$ & proved; proved (written proofs in supplement) & Lem.~\ref{lem:rootedp4}; Thm.~\ref{thm:root}\\
$\ROOT(6)$--$\ROOT(8)$ & computer-assisted & Thm.~\ref{thm:root}\\
$\RAD(8)$--$\RAD(10)$, sharp & computer-assisted & Thm.~\ref{thm:radsharp}\\
REAL-133 for $\lr\le5$ & consequence & Cor.~\ref{cor:five}\\
Band $5<\lr\le6$: two cells left; partial closures & partly computer-assisted & \S\ref{sec:cells}, App.~\ref{app:cells}\\
Bipartite / block classes, all $q$ and $r$ & proved; radius part uses $\mathrm{Struct}(r)$ & \S\ref{sec:classes}\\
Counterexamples to strengthenings & proved (explicit) & \S\ref{sec:neg}\\
Partial TF mass / $\mathrm{EB}_6$ tools (R35--R40) & computer-assisted (referee CWF/CORRECT) & Thm.~G, App.~\ref{app:window}\\
Universal $\mathrm{EB}_6$; unrestricted TF $\GEN_6$; $\GEN_K$ ($K\ge5$ full); $\ROOT(q)$ ($q\ge9$; partial for $q=9$); $\RAD(r)$ ($r\ge11$); TD; $\mathrm{CB}_{\mathrm{high}}$; REAL-133 & open / conjecture & \S\ref{sec:open}\\
\bottomrule
\end{tabular}
\caption{Overview of results and their status.}\label{tab:overview}
\end{table}

\section{Preliminaries}\label{sec:prelim}

For $v\in V(G)$ let $d(v)$ be its degree and $t(v)$ the number of triangles containing $v$. Write $\lambda(v)=\lambda_G(v)=\alpha(G[N(v)])$, $S(G)=\sum_v\lambda(v)$, $e(G)$ for the number of edges and $T(G)$ for the number of triangles.

\begin{lemma}[proved]\label{lem:basic}
Let $G$ be strong $C_4$-free. Then:
\begin{enumerate}
\item two distinct vertices have at most one common neighbour;
\item for every $v$, $G[N(v)]$ is a matching together with isolated vertices, so $\lambda(v)$ equals the number of components of $G[N(v)]$ and $\lambda(v)=d(v)-t(v)$;
\item $S(G)=2e(G)-3T(G)$.
\end{enumerate}
\end{lemma}

\begin{proof}
(1) Two common neighbours $w,w'$ of $x,y$ give the $4$-cycle $xwyw'$. (2) A vertex of $N(v)$ with two neighbours in $N(v)$ would have two common neighbours with $v$. So the components are single vertices and edges, each edge being a triangle at $v$. (3) Sum $d(v)-t(v)$ over $v$.
\end{proof}

The components of $G[N(v)]$ are the \emph{units} of $v$; each is a vertex or an edge. For a vertex set $U$, $G-U$ is the subgraph induced by $V\setminus U$, and $D_G(U)=S(G)-S(G-U)$ is the \emph{deletion loss} of $U$. We write $L(v)=D_G(\{v\})$.

\begin{lemma}[deletion loss; proved]\label{lem:loss}
Let $G$ be strong $C_4$-free and $\emptyset\ne U\subseteq V(G)$. Then
\[
D_G(U)=\sum_{u\in U}\lambda(u)+\sum_{x\notin U}\#\{\text{units of }x\text{ contained in }U\}.
\]
In particular $L(v)=2\lambda(v)-t(v)=2d(v)-3t(v)$.
\end{lemma}

\begin{proof}
For $x\notin U$, deleting $U$ removes exactly the units of $x$ that lie inside $U$; a unit that is an edge with one end in $U$ survives as a vertex. For $U=\{v\}$: the $2t(v)$ neighbours of $v$ in triangles with $v$ keep their value, and each of the other $d(v)-2t(v)$ neighbours loses one. The total is $\lambda(v)+d(v)-2t(v)=2\lambda(v)-t(v)$.
\end{proof}

\begin{lemma}[cut vertices; proved]\label{lem:cut}
Let $v$ be a cut vertex of $G$. Split the components of $G-v$ into two nonempty groups, and let $G_1,G_2$ be the subgraphs induced by the two groups together with $v$. Then $S(G)=S(G_1)+S(G_2)$.
\end{lemma}

\begin{proof}
$G[N(v)]$ has no edges between the groups, so $\lambda_G(v)=\lambda_{G_1}(v)+\lambda_{G_2}(v)$. Every other vertex keeps its neighbourhood.
\end{proof}

We use three published results.

\begin{theorem}[Erd\H{o}s, Saks and S\'os~\cite{ess}]\label{thm:ess}
Every connected graph $G$ has an induced path on at least $2r(G)-1$ vertices.
\end{theorem}

\begin{theorem}[Camby and Schaudt~{\cite[Theorem~5]{cs}}]\label{thm:cs}
Let $k\ge4$ and let $G$ be a connected graph with no induced path on $k$ vertices. Then $G$ has a connected dominating set $X$ such that $G[X]$ has no induced path on $k-2$ vertices or $G[X]$ is isomorphic to the cycle $C_k$.
\end{theorem}

\begin{theorem}[Seinsche~\cite{seinsche}]\label{thm:sein}
If a graph on at least two vertices has no induced path on four vertices, then it or its complement is disconnected.
\end{theorem}

Theorem~\ref{thm:cs} is the implication (i)$\Rightarrow$(ii) of~\cite[Theorem~5]{cs} applied with $H=G$. We write $P_k$ for the path on $k$ vertices, and ``$G$ is $P_k$-free'' means that $G$ has no induced $P_k$, that is, $\pp(G)\le k-1$. We also use the elementary fact that a connected graph of radius $r$ has at least $2r$ vertices (a spanning tree has radius at least $r$, and a tree on $n$ vertices has radius at most $\lfloor n/2\rfloor$).

\begin{lemma}[proved]\label{lem:join}
Let $H$ be a connected strong $C_4$-free $P_4$-free graph on at least two vertices. Then $H$ has a vertex adjacent to all other vertices.
\end{lemma}

\begin{proof}
By Theorem~\ref{thm:sein}, $H$ is the join of two nonempty vertex sets $A$ and $B$. If $|A|,|B|\ge2$, then $a,a'\in A$ and $b,b'\in B$ span the $4$-cycle $aba'b'$. So one side is a single vertex, and it is adjacent to everything.
\end{proof}
\section{Small mass: the range \texorpdfstring{$\lr\le3$}{l<=3}}\label{sec:small}

\begin{lemma}[radius one; proved]\label{lem:r1}
If $G$ is connected and strong $C_4$-free with $r(G)=1$, then $S(G)=2(n-1)-m$, where $m$ is the number of edges of $G[N(u)]$ for a dominating vertex $u$. Hence $\lr(G)<2$, and~\eqref{eq:target} holds.
\end{lemma}

\begin{proof}
Let $v\ne u$. By Lemma~\ref{lem:basic}(1), $v$ has at most one neighbour other than $u$, and such a neighbour lies in $N(u)$. So $G[N(v)]$ is a clique and $\lambda(v)=1$; also $\lambda(u)=(n-1)-m$. Hence $S=2(n-1)-m<2n$ and $\lceil\lr\rceil\le2$. If $G$ is complete then $G\in\{K_2,K_3\}$, $\lr=1$ and $\pp=2=r+1$. Otherwise two nonadjacent vertices of $N(u)$ together with $u$ form an induced $P_3$, so $\pp\ge3\ge r+\lceil\lr\rceil$.
\end{proof}

\begin{lemma}[proved]\label{lem:r2}
If $G$ is connected and strong $C_4$-free with $r(G)\ge2$, then $\pp(G)\ge4$.
\end{lemma}

\begin{proof}
If $G$ were $P_4$-free, Lemma~\ref{lem:join} would give a dominating vertex, so $r(G)\le1$.
\end{proof}

\begin{lemma}[$P_5$-free mass; proved]\label{lem:p5}
If $G$ is connected, strong $C_4$-free and $P_5$-free, then $S(G)\le2n-h$, where $h$ is the number of vertices $v$ with $\lambda(v)\ge3$. In particular $S(G)\le2n$.
\end{lemma}

\begin{proof}
Call $v$ \emph{simplicial} if $\lambda(v)=1$ and \emph{high} if $\lambda(v)\ge3$. For $x\in N(v)$ put $X_x=N(x)\setminus N[v]$; then $N(x)\subseteq\{v\}\cup\{\bar x\}\cup X_x$, where $\bar x$ is the partner of $x$ if $x$ lies in an edge unit of $v$.

\emph{Key fact.} Let $x,y$ lie in different units of $v$. A common vertex of $X_x$ and $X_y$ would give the $4$-cycle $vxwy$, so $X_x\cap X_y=\emptyset$. If $x'\in X_x$ and $y'\in X_y$ are nonadjacent, then $x'xvyy'$ is an induced $P_5$: $x\not\sim y$ because they lie in different units, $x\not\sim y'$ and $y\not\sim x'$ by disjointness, and $v$ is adjacent to neither $x'$ nor $y'$. So $X_x$ is complete to $X_y$. If both are nonempty and $|X_x|\ge2$, a vertex of $X_y$ has two neighbours in $X_x$, which are also neighbours of $x$, a contradiction to Lemma~\ref{lem:basic}(1). Hence if both are nonempty, then $|X_x|=|X_y|=1$.

\emph{Main step.} For each high $v$ we find $\lambda(v)-1$ simplicial vertices whose only high neighbour is $v$, one in each of $\lambda(v)-1$ distinct units of $v$.

Case 1: some $u\in N(v)$ has $|X_u|\ge2$. By the key fact $X_y=\emptyset$ for every $y$ outside the unit of $u$. Then $N(y)\subseteq\{v,\bar y\}$, which is a clique, so $y$ is simplicial; so is $\bar y$. Hence $v$ is the only high neighbour of $y$, and the $\lambda(v)-1$ units other than that of $u$ supply the vertices.

Case 2: $|X_x|\le1$ for all $x\in N(v)$. Then $\lambda(x)\le2$ for $x\in N(v)$. Call a unit \emph{active} if one of its vertices $x$ has $X_x\ne\emptyset$. Suppose two units are active, with $y'\in X_y$ and $z'\in X_z$, and let $x$ lie in a third unit. By the key fact $y'\sim z'$, and $xvyy'z'$ is an induced $P_5$: $x\not\sim y$ (different units), $x\not\sim y',z'$ (disjointness), $v\not\sim y',z'$, and $y\not\sim z'$ (disjointness). As $\lambda(v)\ge3$, a third unit exists, so at most one unit is active. Every vertex of an inactive unit is simplicial and its only high neighbour is $v$. This gives at least $\lambda(v)-1$ units.

\emph{Counting.} The chosen vertices for distinct high vertices are distinct, since each has a unique high neighbour. So $\sum_{v\text{ high}}(\lambda(v)-1)\le\#\{\text{simplicial vertices}\}$, and
\[
S=\sum_v\lambda(v)\le 2n-\#\{\text{simplicial}\}+\sum_{v\text{ high}}(\lambda(v)-2)\le2n-h.\qedhere
\]
\end{proof}

\begin{lemma}[$C_6$ attachments; proved]\label{lem:c6}
Let $G$ be connected, strong $C_4$-free and $P_6$-free with a dominating induced cycle $C=v_0v_1\cdots v_5$, and let $W=V(G)\setminus V(C)$. Then every $w\in W$ is adjacent to exactly one antipodal pair $\{v_i,v_{i+3}\}$ and to no other vertex of $C$, distinct vertices of $W$ use distinct pairs, and $W$ is independent. Hence $|W|\le3$, $S(G)=12+4|W|<3n$, and $r(G)\le2$ if $W\ne\emptyset$.
\end{lemma}

\begin{proof}
Indices are modulo $6$. Let $w\in W$; it has a neighbour on $C$. If $w$ has exactly one neighbour $v_i$, then $wv_iv_{i+1}v_{i+2}v_{i+3}v_{i+4}$ is an induced $P_6$. If $N(w)\cap C=\{v_i,v_{i+1}\}$, then $wv_{i+1}v_{i+2}v_{i+3}v_{i+4}v_{i+5}$ is an induced $P_6$. Two neighbours at distance $2$ on $C$ give a $4$-cycle. A set of at least three vertices of $C_6$ contains two vertices at distance $2$. So $N(w)\cap C$ is an antipodal pair. Two vertices of $W$ with the same pair give two common neighbours. If $w\sim w'$ with distinct pairs, some $v_i\in N(w)$ and $v_j\in N(w')$ are adjacent on $C$, and $wv_iv_jw'$ is a $4$-cycle. There are no triangles, $e(G)=6+2|W|$, and Lemma~\ref{lem:basic}(3) gives $S=12+4|W|$. Finally, if $w$ uses $\{v_i,v_{i+3}\}$, then every vertex of $C$ is within distance $2$ of $v_i$ (the vertex $v_{i+3}$ through $w$), and each $w'\in W$ has a neighbour in $\{v_{i-1},v_i,v_{i+1}\}$, because each of the three antipodal pairs meets this set. So $v_i$ has eccentricity $2$.
\end{proof}

The included checker (Section~\ref{sec:data}) also verifies Lemma~\ref{lem:c6} by brute force: all $63$ traces, all $63\cdot63\cdot2$ configurations of two outside vertices, and all type subsets.

\begin{lemma}[proved]\label{lem:ecc2}
Let $G$ be connected, strong $C_4$-free and $P_6$-free. Then either $G$ has a vertex of eccentricity at most $2$, or $G=C_6$.
\end{lemma}

\begin{proof}
Apply Theorem~\ref{thm:cs} with $k=6$ to get a connected dominating set $X$. If $G[X]\cong C_6$, Lemma~\ref{lem:c6} gives $G=C_6$ or a vertex of eccentricity $2$. Otherwise $G[X]$ is $P_4$-free. If $|X|=1$, the vertex in $X$ dominates $G$. If $|X|\ge2$, Lemma~\ref{lem:join} gives $u\in X$ adjacent to all of $X\setminus\{u\}$, and since $X$ dominates $G$, every vertex is within distance $2$ of $u$.
\end{proof}

\begin{lemma}[$P_6$-free mass; proved]\label{lem:p6}
If $G$ is connected, strong $C_4$-free and $P_6$-free, then $S(G)\le3n$, with equality if and only if $G$ is the Petersen graph.
\end{lemma}

\begin{proof}
If $G=C_6$ then $S=12<18$. Otherwise, by Lemma~\ref{lem:ecc2}, fix $u$ of eccentricity at most $2$. Put $N_1=N(u)$, $N_2=V\setminus N[u]$, and let $m_1$ be the number of edges of $G[N_1]$. By Lemma~\ref{lem:basic}(1) each $c\in N_2$ has exactly one neighbour in $N_1$, its \emph{parent}. For $a\in N_1$ let $C_a$ be the set of its children. An edge $cd$ of $G[N_2]$ with distinct parents is a \emph{cross edge}. Let $E'$ be the set of cross edges, and let $D_c$ be the set of cross neighbours of $c$.

\emph{Local values.} $\lambda(u)=|N_1|-m_1$. For $a\in N_1$ with partner $\bar a$ (if any), no child of $a$ is adjacent to $u$ or $\bar a$, so the units of $a$ are $\{u,\bar a\}$ (or $\{u\}$) together with the components of $G[C_a]$; hence $\lambda(a)=1+\alpha(G[C_a])$. For $c\in C_a$: $c$ has at most one neighbour $s\in C_a$ (two would be common neighbours of $a$ and $c$), the vertices of $D_c$ are not adjacent to $a$, and $s$ is adjacent to no $d\in D_c$ (else $c,a,s,d$ is a $4$-cycle). So $\lambda(c)=1+\alpha(G[D_c])$. Using $\sum_a|C_a|=|N_2|$ and $n=1+|N_1|+|N_2|$ we get
\begin{equation}\label{eq:p6id}
3n-S=3+|N_1|+m_1+\sum_a\bigl(|C_a|-\alpha(G[C_a])\bigr)+|N_2|-\sum_{c}\alpha(G[D_c])\ \ge\ 3+|N_1|+|N_2|-2|E'|.
\end{equation}
So it suffices to show $2|E'|\le3+|N_1|+|N_2|$.

\emph{Fact 1.} Let $cd\in E'$ with $c\in C_y$ and $d\in C_z$. Then $y\not\sim z$, else $ycdz$ is a $4$-cycle. Let $y'\in N_1\setminus\{y,z\}$ and $x\in C_{y'}$. Then $x$ is adjacent to $c$ or $d$. Indeed, otherwise: if $y'\not\sim y$, then $xy'uycd$ is an induced $P_6$; if $y'\sim y$, then $y'\not\sim z$ (since $G[N_1]$ is a matching) and $xy'uzdc$ is an induced $P_6$. Two children of $y'$ adjacent to the same vertex of $\{c,d\}$ would be two common neighbours of $y'$ and that vertex. So $|C_{y'}|\le2$.

\emph{Fact 2.} If $E'\ne\emptyset$, at most three vertices of $N_1$ have children. Suppose $cd\in E'$ has parents $y,z$, and $y_1,y_2\in N_1\setminus\{y,z\}$ have children $x_1,x_2$. Let $H=G[\{c,d,x_1,x_2\}]$. These four vertices have distinct parents, so every edge of $H$ is a cross edge, and by Fact~1 each of the other two vertices is adjacent to an endpoint of that edge. Hence $H$ has no induced $K_2\cup K_1$. So the complement of $H$ has no induced $P_3$, that is, it is a disjoint union of cliques, and $H$ is complete multipartite. A complete multipartite graph on four vertices with at least one edge and no $4$-cycle subgraph is $K_{1,3}$. Let its centre be $c^*$ and its leaves $a,b,d^*$, with parents $a_0,b_0,d_0$. Since $G[N_1]$ is a matching, two of these parents are nonadjacent, say $a_0\not\sim b_0$. Then $a_0,u,b_0,b,c^*,d^*$ is an induced $P_6$: the leaves are pairwise nonadjacent, and each child is adjacent only to its own parent in $N_1$. This is a contradiction. (The included checker also confirms Fact~2 by exhausting the $2^{10}$ possible adjacency patterns of the nine vertices $u,y,z,y_1,y_2,c,d,x_1,x_2$.)

\emph{Counting.} A child has at most one cross neighbour in each other class $C_{y'}$ (two would be common neighbours with $y'$), so its number of cross neighbours is at most $k-1$, where $k\le3$ is the number of vertices of $N_1$ with children. If $E'=\emptyset$ we are done. If $k=2$, every child has at most one cross neighbour, and $2|E'|\le|N_2|$. Let $k=3$, with classes $C_1,C_2,C_3$. If a child in $C_i$ has two cross neighbours, in $C_j$ and $C_l$, then Fact~1 applied to these two cross edges gives $|C_j|,|C_l|\le2$. Distinct such children of $C_i$ have distinct neighbours in $C_j$ (else two common neighbours with $y_i$), so $C_i$ contains at most two children with two cross neighbours, and there are at most six in total. Hence $2|E'|\le|N_2|+6\le|N_2|+3+|N_1|$, because $|N_1|\ge k=3$.

\emph{Equality.} The Petersen graph is $P_6$-free: if $v_0\cdots v_5$ were an induced path, then $v_0$ would have a common neighbour with each of $v_3,v_4,v_5$ (diameter $2$); these lie off the path (the only path neighbour of $v_0$ is $v_1$, which is not adjacent to them) and are distinct (girth $5$), but $v_0$ has only two neighbours off the path. Also $\lambda\equiv3$, so $S=30=3n$. Conversely, suppose $S=3n$. Then $G\ne C_6$; the dominating-$C_6$ case is excluded because there $S<3n$ (Lemma~\ref{lem:c6}); and $E'\ne\emptyset$, because otherwise \eqref{eq:p6id} gives $3n-S\ge3+|N_1|+|N_2|>0$. If $k=2$ then $3n-S\ge3+|N_1|>0$. So $k=3$, and equality throughout forces: $|N_1|=3$; $m_1=0$; no edges inside any $C_a$; every $D_c$ is independent; and exactly six children have two cross neighbours, two in each class. By Fact~1 every class has at most two children, so every class has exactly two, and every child has exactly two cross neighbours. The cross edges then form a $2$-regular graph on the six children with no triangle (each $D_c$ is independent), that is, a $6$-cycle. Two children of the same parent at distance $2$ on this cycle would have two common neighbours, and they are not adjacent, so they are antipodal. This determines $G$ up to isomorphism, and the Petersen graph has exactly this structure: in it, the six vertices at distance $2$ from a vertex induce a $6$-cycle.
\end{proof}

\begin{theorem}[Theorem A; proved]\label{thm:A}
Let $G$ be connected and strong $C_4$-free with $n\ge2$. If $\lr(G)\le3$, or if $r(G)\le2$ and $\lr(G)\le4$, then $\pp(G)\ge r(G)+\lr(G)$.
\end{theorem}

\begin{proof}
Write $r=r(G)$ and $c=\lceil\lr(G)\rceil$; we need $\pp\ge r+c$.
\begin{itemize}
\item $r=1$: Lemma~\ref{lem:r1}.
\item $r=2$, $c\le2$: Lemma~\ref{lem:r2}.
\item $r=2$, $c=3$: if $\pp\le4$, then Lemma~\ref{lem:p5} gives $\lr\le2$, a contradiction. So $\pp\ge5$.
\item $r=2$, $c=4$: if $\pp\le5$, then Lemma~\ref{lem:p6} gives $\lr\le3$, a contradiction. So $\pp\ge6$.
\item $r=3$, $c\le2$: Theorem~\ref{thm:ess} gives $\pp\ge5$.
\item $r=3$, $c=3$: if $\pp\le5$, then by Lemma~\ref{lem:ecc2} and $r=3$ we have $G=C_6$, with $\lr=2$, a contradiction. So $\pp\ge6$.
\item $r\ge4$, $c\le3$: Theorem~\ref{thm:ess} gives $\pp\ge2r-1\ge r+3$. \qedhere
\end{itemize}
\end{proof}

Equality occurs, for instance, for $C_5$ ($r=2$, $\lr=2$, $\pp=4$), the Petersen graph ($r=2$, $\lr=3$, $\pp=5$) and $C_6$ ($r=3$, $\lr=2$, $\pp=5$). We do not characterise the equality cases.

\section{Reduction to bounded induced paths}\label{sec:band}

\begin{theorem}[Theorem B; proved]\label{thm:B}
Let $G$ be a connected strong $C_4$-free graph that violates Conjecture~\ref{conj:133}. Then $\lr(G)>3$. Put $K=\lceil\lr(G)\rceil-1\ge3$. Then either
\begin{enumerate}
\item[(a)] $r(G)\le K$ and $\pp(G)\le r(G)+K\le2K$, so $G$ violates $\GEN_K$; or
\item[(b)] $r(G)=K+1\ge4$ and $\pp(G)=2r(G)-1$, and $\lr(G)>r(G)-1$, so $G$ violates $\RAD(K+1)$.
\end{enumerate}
\end{theorem}

\begin{proof}
Write $r=r(G)$, $\pp=\pp(G)$ and $c=\lceil\lr(G)\rceil=K+1$. Then $\lr>3$ by Theorem~\ref{thm:A}. A violation means $\pp\le r+c-1=r+K$. Theorem~\ref{thm:ess} gives $2r-1\le\pp\le r+K$, so $r\le K+1$. If $r\le K$, then $\pp\le r+K\le2K$, and $S>Kn$ because $\lr>K$. If $r=K+1$, then $2r-1\le\pp\le 2K+1=2r-1$, and $\lr>K=r-1$.
\end{proof}

Since $\GEN_K$ ignores the radius, it is stronger than what alternative~(a) needs: a counterexample in~(a) also satisfies $\pp\le r+K$. We know of no strong $C_4$-free graph violating any $\GEN_K$. Note, however, that $\GEN_K$ for large $K$ might fail even if the conjecture holds; see Remark~\ref{rem:moore}.

\begin{proposition}[proved]\label{prop:theta}
For $k\ge3$ let $\theta_k$ consist of two vertices $x,y$ joined by $k$ internally disjoint paths $xa_ib_iy$. Then $\theta_k$ is strong $C_4$-free, $n=2k+2$, $r=3$, $\pp=6$ and $S=6k=3n-6$. Thus $\pp-r-\lr=6/n\to0$.
\end{proposition}

\begin{proof}
Every cycle consists of two complete branches and so has length $6$; in particular there is no $4$-cycle. There are no triangles, so $S=2e=6k$. Every vertex has eccentricity $3$: $d(x,y)=3$ and $d(a_i,b_j)=3$ for $i\ne j$. The path $a_jxa_ib_iyb_l$ (with $j,i,l$ distinct) is induced. Conversely, let $P$ be an induced path. If $P$ avoids $x$ and $y$, it lies inside one edge $a_ib_i$. If $P$ contains $x$ but not $y$, then $x$ has at most two neighbours on $P$, each followed on $P$ by at most its $b$-vertex, so $|P|\le5$. If $P$ contains $x$ and $y$, then every $a_j$ on $P$ is adjacent on $P$ to $x$ and every $b_j$ to $y$, at most one branch lies entirely on $P$, and so $P$ consists of $x$, $y$, at most one full branch and at most one further neighbour of each of $x$ and $y$. Hence $|P|\le6$.
\end{proof}

So $\GEN_3$ is sharp up to $O(1/n)$ already among graphs of radius~$3$, and no argument with a uniform positive margin in $\lr$ can prove it.

\section{The case \texorpdfstring{$K=3$}{K=3}: proof of \texorpdfstring{$\GEN_3$}{GEN3}}\label{sec:gen3}

\begin{theorem}[$\GEN_3$; proved, with three finite certificates; formalised]\label{thm:gen3}
Every strong $C_4$-free graph $G$ with no induced $P_7$ satisfies $S(G)\le3n$. Connectivity is not needed.
\end{theorem}

Combined with Theorem~\ref{thm:B}, Theorem~\ref{thm:gen3} settles REAL-133 for every graph with $\pp\le6$, and in particular closes the cell $r=3$, $3<\lr\le4$. The proof is a minimum-counterexample argument. Section~\ref{sec:lean} describes its Lean formalisation.

\subsection{Minimum counterexamples}
\begin{lemma}[no rooted $P_4$; proved]\label{lem:rootedp4}
Let $F$ be connected and strong $C_4$-free with a vertex $s$ such that no induced $P_4$ of $F$ has $s$ as an end. Then $S(F)\le2(|F|-1)-m$, where $m$ is the number of edges of $F[N(s)]$. This is $\ROOT(3)$.
\end{lemma}

\begin{proof}
A vertex at distance $3$ from $s$ would end a geodesic, hence an induced $P_4$, starting at $s$. So $V=\{s\}\cup N_1\cup N_2$ with $N_i$ the $i$-th distance layer. Each $y\in N_2$ has a unique parent $a\in N_1$. A neighbour $z\in N_2$ of $y$ with a different parent would give the induced path $sayz$. So $y$ has at most one neighbour in $N_2$, a sibling (two siblings would be common neighbours of $a$ and $y$), and $N(y)$ is a clique; thus $\lambda(y)=1$. For $a\in N_1$ we have $N(a)\subseteq\{s,\bar a\}\cup C_a$, where $C_a$ is the set of children of $a$ and $\bar a$ is the neighbour of $a$ in $N_1$, if any. Since $s$ and $\bar a$ lie in one unit, $\lambda(a)\le1+|C_a|$. Finally $\lambda(s)=|N_1|-m$. Summing, $S\le(|N_1|-m)+(|N_1|+|N_2|)+|N_2|=2(|F|-1)-m$.
\end{proof}

\begin{theorem}[proved]\label{thm:red}
Let $G$ be a counterexample to $\mathrm{GEN}_3$ of minimum order, that is, $G$ is connected and strong $C_4$-free, $\pp(G)\le6$, $S(G)\ge3n+1$, and $n$ is minimum. Then:
\begin{enumerate}
\item[(a)] $G$ is $2$-connected;
\item[(b)] $S(G)-S(G-U)\ge3|U|+1$ for every nonempty proper $U\subset V(G)$;
\item[(c)] $2\lambda(v)-t(v)\ge4$ for every $v$; in particular $\lambda(v)\ge2$;
\item[(d)] every vertex with $\lambda(v)=2$ has degree $2$ and nonadjacent neighbours, and no two vertices of degree $2$ are adjacent.
\end{enumerate}
\end{theorem}

\begin{proof}
(a) Since $S(G)\le 2<3n$ for $n\le2$, we have $n\ge3$, so it suffices to exclude cut vertices. Let $v$ be a cut vertex. For each component $C$ of $G-v$ let $\rho_C$ be the largest order of an induced path in $G[C\cup\{v\}]$ with end $v$. Two such paths from different components meet only in $v$ and have no edges between them, so together they form an induced path on $\rho_C+\rho_{C'}-1\le6$ vertices. So some component $C_1$ has $\rho_{C_1}\le3$. Put $G_1=G[C_1\cup\{v\}]$ and $G_2=G-C_1$, with $n_i=|G_i|$, so $n=n_1+n_2-1$ and $n_1\ge2$. Both graphs are connected, smaller than $G$, and $P_7$-free. Lemma~\ref{lem:rootedp4} gives $S(G_1)\le2(n_1-1)$, and minimality gives $S(G_2)\le3n_2$. By Lemma~\ref{lem:cut}, $S(G)\le2(n_1-1)+3n_2=3n-(n_1-1)<3n$, a contradiction.

(b) Every component $H$ of $G-U$ is connected, smaller, and $P_7$-free, so $S(H)\le3|H|$ (for $|H|=1$, $S(H)=0$). Summing, $S(G-U)\le3(n-|U|)$, and $S(G)\ge3n+1$ gives the claim.

(c) Take $U=\{v\}$ in (b) and use Lemma~\ref{lem:loss}. Since $t(v)\ge0$, $2\lambda(v)\ge4$.

(d) If $\lambda(v)=2$, then (c) forces $t(v)=0$, so $d(v)=2$ and the neighbours are nonadjacent. Every vertex of degree $2$ has $\lambda=2$ by (c), hence $t=0$. Let $u\sim w$ both have degree $2$, with other neighbours $x$ and $y$. Then $x\ne y$, since $t(u)=0$; $\{u\}$ is a unit of $x$ and $\{w\}$ a unit of $y$. Lemma~\ref{lem:loss} with $U=\{u,w\}$ gives a loss of $2+2+1+1=6<7$, contradicting (b).
\end{proof}

The disconnected case follows from the connected one by summing over components (an isolated vertex has $\lambda=0$). Call a graph \emph{admissible} if it is strong $C_4$-free and $P_7$-free, $L(v)=2\lambda(v)-t(v)\ge4$ for every vertex $v$, and no two adjacent vertices have degree~$2$. By Theorem~\ref{thm:red}(b)--(d), a counterexample to $\GEN_3$ of minimum order is admissible. In an admissible graph $\lambda\ge2$ everywhere; $\lambda(v)=2$ forces $t(v)=0$ and $d(v)=2$, so every degree-$2$ vertex has nonadjacent neighbours and lies in no triangle; and every vertex of degree at least $3$, a \emph{branch vertex}, has $\lambda\ge3$. Let $B$ be the set of branch vertices, $Z$ the set of degree-$2$ vertices (\emph{ears}), $D=G[B]$, $d_D(v)=|N(v)\cap B|$ and $q(v)=|N(v)\cap Z|$. All triangles lie in $D$, so $t_D(v)=t(v)$.

\begin{proposition}[proved, given Lemmas~\ref{lem:ears}--\ref{lem:mixed}]\label{prop:adm}
Every admissible graph satisfies $S\le3n$. Hence Theorem~\ref{thm:gen3} holds.
\end{proposition}

\begin{proof}
Ears have nonadjacent neighbours and $Z$ is independent, so each ear contributes $\lambda=2$ and one singleton unit to each of its two neighbours: $S(G)=S(D)+4|Z|$. Since $|Z|=\frac12\sum_{v\in B}q(v)$ and $n=|B|+|Z|$,
\[
S-3n=\sum_{v\in B}\Bigl(d_D(v)-t_D(v)-3+\tfrac12q(v)\Bigr).
\]
By Lemmas~\ref{lem:ears}, \ref{lem:branch} and~\ref{lem:mixed} every summand is at most $0$: for $(d_D,t_D)=(0,0)$ we have $q\le4$ and the summand is at most $-1$; for $(1,0)$ it is at most $0$; for $(2,0)$ the two branch neighbours are nonadjacent, so $q\le2$ and the summand is at most $0$; for $(2,1)$, $q\le4$ gives at most $0$; for $(3,t)$, $q=0$ gives $-t\le0$. Lemma~\ref{lem:branch} excludes $d_D\ge4$.
\end{proof}

\subsection{Local lemmas}

\begin{lemma}[at most four ears; proved]\label{lem:ears}
In an admissible graph $q(v)\le4$ for every $v$.
\end{lemma}

\begin{proof}
Let $u_i$ be ears at $v$ with other neighbours $a_i$. The $a_i$ are distinct (else $v u_i a u_j$ is a $C_4$), nonadjacent to $v$ (ears lie in no triangle) and of degree at least $3$ (no adjacent degree-$2$ pair). Put $E_i=N(a_i)\setminus\{u_i\}$; then $|E_i|\ge2$, and $E_i\cap N[v]=\emptyset$, since a second common neighbour of $v$ and $a_i$ would give a $C_4$.

\emph{Claim.} If $a_i\not\sim a_j$, then $E_i=\{w,p\}$, $E_j=\{w,q\}$ with $pq\in E$, $w\not\sim p$ and $w\not\sim q$; hence $a_i,a_j$ have degree $3$ and lie in no triangle. Indeed $|E_i\cap E_j|\le1$; for $p\in E_i\setminus E_j$ and $q\in E_j\setminus E_i$ with $p\not\sim q$ the path $p a_i u_i v u_j a_j q$ is induced (all chords are excluded by $d(u_i)=d(u_j)=2$, $E_i\cap N[v]=E_j\cap N[v]=\emptyset$, $a_i\not\sim a_j$ and the choice of $p,q$). So the private parts are complete to each other. Both are nonempty, and if one had two vertices $p,p'$, then $a_i p q p'$ would be a $C_4$; so both have size one and $|E_i\cap E_j|=1$. An edge $wp$ would give the $C_4$ $w a_j q p$.

\emph{No three pairwise nonadjacent $a_i$.} Otherwise the sets $E_i$ pairwise meet in exactly one element. If the three intersections are the three pairs of a $3$-set $\{\alpha,\beta,\gamma\}$, the claim forces the triangle $\alpha\beta\gamma$, and $a_1\sim\alpha,\beta$ gives the $C_4$ $a_1\alpha\gamma\beta$. Otherwise $E_i=\{w,p_i\}$ with $p_1p_2p_3$ a triangle and $w\not\sim p_i$, and $v u_1 a_1 w a_2 p_2 p_3$ is an induced $P_7$.

\emph{Five ears.} Then $F=G[\{a_1,\dots,a_5\}]$ has $\alpha(F)\le2$. If $F$ has a vertex $a$ adjacent to the other four, these induce a matching (Lemma~\ref{lem:basic}), and since $\alpha(F)\le2$ it is a perfect matching $\{b_1b_2,c_1c_2\}$; $b_1\not\sim c_1$, so by the claim $b_1$ lies in no triangle, but $ab_1b_2$ is one. Otherwise, by the claim every vertex of $F$ that has a non-neighbour in $F$ has degree $3$ in $G$ and lies in no triangle; so $F$ is triangle-free with maximum degree at most $2$, $C_4$-free and has $\alpha(F)\le2$, hence $F=C_5=a_0\cdots a_4$, and $u_0 v u_1 a_1 a_2 a_3 a_4$ is an induced $P_7$.
\end{proof}

\paragraph{Attachment certificates.} Let $H$ be an induced subgraph of an admissible graph $G$, let $R\subseteq V(H)$ be a set of vertices known to have all their $G$-neighbours in $H$, and $K\subseteq V(H)$ a set of vertices known to be branch vertices. A \emph{kept trace} is a nonempty $A\subseteq V(H)\setminus R$ such that $H$ plus one new vertex with neighbourhood $A$ is strong $C_4$-free and $P_7$-free. Call $x\in V(H)$ \emph{raisable} if some kept trace contains $x$ and misses $N_H(x)$. If $x$ is not raisable, every neighbour $y$ of $x$ outside $H$ has a kept trace $N(y)\cap V(H)$ (since $H+y$ is induced in $G$) containing $x$ and some old neighbour $z$ of $x$; so $xyz$ is a new triangle, $y$ joins the previously singleton unit of $z$, and $\lambda_G(x)=\lambda_H(x)$ and $L_G(x)\le L_H(x)-(\text{number of new neighbours})$. Hence $H$ cannot occur in an admissible graph if it has (i) a nonraisable $x$ with $L_H(x)<4$, or (ii) a nonraisable $x\in K$ with $\lambda_H(x)<3$, or (iii) an edge $xy$ with $d_H(x)=d_H(y)=2$, both ends nonraisable and $L_H(x)=L_H(y)=4$ (they would remain adjacent degree-$2$ vertices of $G$). No bound on $|V(G)|$ is involved.

\begin{lemma}[two separated forks; finite classification]\label{lem:forks}
An admissible graph contains no induced $7$-vertex tree $T$ with edges $va,vb,aa_1,aa_2,bb_1,bb_2$.
\end{lemma}

\begin{proof}[Proof outline]
Since $\lambda(a_i)\ge2$ there is $A_i\in N(a_i)\setminus N[a]$, and similarly $B_j\in N(b_j)\setminus N[b]$. These lie outside $T$, are nonadjacent to $v$ (a $C_4$ with $a$ or $b$) and are pairwise distinct: on the same side by $C_4$-freeness, and $A_i=B_j$ by the induced path $a_{3-i}aa_iA_ib_jbb_{3-j}$. Call $A_i$ \emph{targeted} if $A_i\sim b$ (and $B_j$ if $B_j\sim a$). A case analysis with explicit induced $P_7$'s shows that all four witnesses are targeted. Then $a,b$ are nonadjacent hubs with common neighbour $v$ and four internally disjoint paths $ax_iy_ib$ ($X=\{a_1,a_2,B_1,B_2\}\subseteq N(a)$, $Y=\{A_1,A_2,b_1,b_2\}\subseteq N(b)$). On these $11$ vertices the only undetermined pairs are matching edges inside $X$ and inside $Y$. Of the $10\times10=100$ configurations, $79$ are strong $C_4$-free, $43$ are also $P_7$-free, and each of the $43$ is excluded by an attachment certificate with $R=K=\emptyset$ ($2047$ candidate traces each, of which $5$ to $13$ are kept).
\end{proof}

\begin{lemma}[at most three branch neighbours; finite classification]\label{lem:branch}
In an admissible graph $d_D(v)\le3$ for every $v$.
\end{lemma}

\begin{proof}[Proof outline]
Let $u_1,\dots,u_4$ be branch neighbours of $v$. As $\lambda(u_i)\ge3$, pick $x_i,y_i$ in two distinct units of $u_i$ other than the unit of $v$; they are nonadjacent and lie outside $N[v]$. These child classes are disjoint, no child is adjacent to another selected parent, and between two classes the edges form a matching. The parents induce a matching. Adjacent parents have anticomplete classes ($C_4$), and nonadjacent parents have at least one edge between their classes, since otherwise $v,u_i,u_j$ and the four children induce the tree of Lemma~\ref{lem:forks}. So each nonadjacent pair of parents has $6$ options. Of the $6^6=46656$, $6^5=7776$ and $6^4=1296$ configurations (no, one, two parent edges), $24688$, $6288$ and $1048$ are strong $C_4$-free and none is $P_7$-free.
\end{proof}

\begin{lemma}[mixed neighbourhoods; finite classification]\label{lem:mixed}
In an admissible graph: \textup{(a)} $d_D(v)=3$ implies $q(v)=0$; \textup{(b)} if $d_D(v)=2$ and the two branch neighbours are nonadjacent, then $q(v)\le2$.
\end{lemma}

\begin{proof}[Proof outline]
Select three branch neighbours and one ear for (a), and two nonadjacent branch neighbours and three ears for (b). Branch parents get two children as in Lemma~\ref{lem:branch}; an ear parent has a unique child (its other neighbour), a branch vertex nonadjacent to $v$ and to the other parents and distinct from the other children. Ear--branch class pairs have $3$ options, ear--ear pairs $2$, nonadjacent branch--branch pairs $6$ (Lemma~\ref{lem:forks}), and adjacent branch parents are anticomplete. Use $R=$ the ear parents and $K=\{v\}\cup{}$branch parents${}\cup{}$ear children. For (a) with independent parents there are $5832$ configurations, $4320$ strong $C_4$-free and $32$ $P_7$-free, all excluded by attachment certificates; with one parent edge, $972$, $876$ and $0$. For (b), $34992$, $17088$ and $18$, all $18$ excluded.
\end{proof}

The finite parts of Lemmas~\ref{lem:forks}--\ref{lem:mixed} were reproduced by three independent implementations (Section~\ref{sec:method}). As a sanity check (not used in the proof), among all connected strong $C_4$-free $P_7$-free graphs with $n\le14$ the only admissible graphs are four graphs on $8$, $9$, $10$ and $11$ vertices, one of them the Petersen graph. Earlier versions of this paper reduced $\GEN_3$ to a search over ``skeletons'' of branch vertices and verified it for $n\le20$; that material is superseded by Theorem~\ref{thm:gen3} and has been removed.

\subsection{Lean formalisation}\label{sec:lean}
Theorem~\ref{thm:gen3} is formalised in Lean~4 (v4.35.0) with Mathlib~\cite{lean4,mathlib}. The main theorem reads
\begin{center}
\texttt{theorem gen3 (hC : StrongC4Free G) (hP : P7Free G) : S G \(\le\) 3 * Fintype.card V}
\end{center}
for an arbitrary finite simple graph \texttt{G} on a vertex type \texttt{V}; here \texttt{StrongC4Free} says that two distinct vertices have at most one common neighbour, \texttt{P7Free} says that no injective map $\mathrm{Fin}\,7\to V$ is an induced path, and \texttt{S G} is the sum over $v$ of the independence number of the induced subgraph on $N(v)$ (Mathlib's \texttt{indepNum}). The formalisation follows the proof above with one simplification: a single generic configuration checker, with a machine-checked soundness theorem, verifies case-split trees whose leaves are $C_4$ witnesses, induced-$P_7$ witnesses, degree witnesses, contradictions, or iterated attachment extensions. With it, Lemmas~\ref{lem:ears}, \ref{lem:branch}, \ref{lem:mixed}(a) and \ref{lem:mixed}(b) become small hand-written setups plus one certificate each ($1008$, $14969$, $6419$ and $3828$ nodes); Lemma~\ref{lem:forks} is not needed. All finite checks run in the kernel by \texttt{decide +kernel} (no \texttt{native\_decide}). The axioms of \texttt{gen3} are \texttt{propext}, \texttt{Classical.choice} and \texttt{Quot.sound}; the certificate checks use only \texttt{propext}. A clean build of the development (with Mathlib cached) takes about three minutes.

\section{The case \texorpdfstring{$K=4$}{K=4}: \texorpdfstring{$\GEN_4$}{GEN4}}\label{sec:gen4}

\begin{theorem}[$\GEN_4$; computer-assisted]\label{thm:gen4}
Every connected strong $C_4$-free graph with no induced $P_9$ satisfies $S(G)\le4n$.
\end{theorem}

\begin{proof}[Proof outline]
(i) Let $G$ be a counterexample of minimum order, so $\pp(G)\le8$ and $S(G)\ge4n+1$. Every component of $G-U$ is smaller, so $D_G(U)\ge4|U|+1$ for every nonempty proper $U$; in particular $L(v)\ge5$ and $\lambda(v)\ge3$ for every $v$.

(ii) \emph{Attachment lemma.} If $H$ is an induced subgraph of $G$ and $x\in V(H)$ has $L_H(x)<5\le L_G(x)$, then some $y\notin V(H)$ adjacent to $x$ has $N(y)\cap N_H(x)=\emptyset$. Otherwise every outside neighbour of $x$ joins a distinct old singleton unit of $x$, so $\lambda_G(x)=\lambda_H(x)$ and $L_G(x)\le L_H(x)$. The trace $A=N(y)\cap V(H)$ contains $x$, misses $N_H(x)$, and $H+y$ is strong $C_4$-free and $P_9$-free.

(iii) \emph{Finite lemma.} No strong $C_4$-free $P_9$-free graph with $L\ge5$ everywhere contains an induced $K_{1,5}$. The certificate is a finite rooted directed acyclic graph of records. The root is $K_{1,5}$. Each record $H$ names a vertex $x$ with $L_H(x)<5$ and lists, for \emph{every} trace $A$ as in (ii), a child record isomorphic to $H+y$, with an explicit isomorphism. There are $3796$ records ($2287$ terminal, that is, without any admissible trace), $6679$ arcs and $1{,}929{,}045$ examined $C_4$-safe traces; the largest record has $27$ vertices. If $G$ contained $K_{1,5}$, following (ii) through the certificate, carrying an induced embedding into $G$, would reach a terminal record, a contradiction. Every arc raises the order by one, so the walk terminates; no bound on $|G|$ is assumed.

(iv) A vertex with $\lambda\ge5$ yields an induced $K_{1,5}$ (one vertex per unit). So $3\le\lambda\le4$ everywhere and $S-4n=-\#\{v:\lambda(v)=3\}\le0$, a contradiction.
\end{proof}

The bound $\lambda\le4$ in (iv) is attained: $C_8$ with four vertices $a_i\sim c_i,c_{i+4}$ and a root adjacent to all $a_i$ has $13$ vertices, $\pp=7$ and $\lambda=4$ at the root. With Theorem~\ref{thm:B}, $\GEN_4$ also shows that $\lr>4$ implies $\pp\ge9$.

\section{The radius-extremal side: \texorpdfstring{$\RAD(r)$}{RAD(r)} for \texorpdfstring{$r\le7$}{r<=7}}\label{sec:rad}
\subsection{Radius four}\label{sec:r4}
A finite classification settles the smallest case $\RAD(4)$.

\begin{lemma}[finite classification]\label{lem:c8}
Let $G$ be connected, strong $C_4$-free and $P_8$-free, with a dominating induced $8$-cycle. If $r(G)\ge4$, then $G=C_8$.
\end{lemma}

\begin{proof}
Let $C$ be the cycle and $W=V(G)\setminus V(C)$. Of the $255$ nonempty traces on $C$, exactly $32$ give a strong $C_4$-free, $P_8$-free graph $C+w$, and each of them has at least two vertices. So two vertices of $W$ with the same trace would have two common neighbours, and the vertices of $W$ have distinct types. List them in increasing order of type; every prefix induces a strong $C_4$-free, $P_8$-free subgraph, since both properties are hereditary. Adding vertices of increasing type, with every adjacency to earlier vertices and pruning at a $4$-cycle or an induced $P_8$, produces $790$ such configurations ($73$ isomorphism classes), with at most $13$ vertices; the bound $13$ is an outcome of the search, not an input. Among them only $C_8$ itself has radius at least $4$.
\end{proof}

\begin{theorem}[radius four; proved, with a finite classification]\label{thm:r4}
Let $G$ be connected and strong $C_4$-free with $r(G)=4$ and $\pp(G)=7$. Then $S(G)\le2n$, i.e.\ $\lr(G)\le2$.
\end{theorem}

\begin{proof}
By Theorem~\ref{thm:cs} with $k=8$, $G$ has a connected dominating set $X$ with $G[X]$ $P_6$-free or $G[X]\cong C_8$. In the second case Lemma~\ref{lem:c8} gives $G=C_8$ and $\lr=2$. In the first case, by Lemma~\ref{lem:ecc2} applied to $G[X]$, either $G[X]$ has a vertex of eccentricity at most $2$ in $G[X]$, which then has eccentricity at most $3$ in $G$ since $X$ dominates $G$, contradicting $r(G)=4$; or $G[X]=C_6$. So let $C=v_0\cdots v_5$ (indices modulo $6$) be a dominating induced $6$-cycle and $W=V(G)\setminus V(C)$.

\emph{Traces.} As in the proof of Lemma~\ref{lem:c6}, strong $C_4$-freeness allows only three kinds of trace $N(w)\cap C$: a single vertex, a consecutive pair $\{v_i,v_{i+1}\}$, or an antipodal pair. If some $w$ has trace $\{v_i,v_{i+3}\}$, then $v_i$ is within distance $2$ of every vertex of $C$, hence within distance $3$ of every vertex of $G$, a contradiction. So there are no antipodal traces. Let $C_i$ be the set of vertices with trace $\{v_i\}$.

\emph{Witnesses.} Fix $i$. Since $r(G)=4$, some vertex is at distance at least $4$ from $v_{i+3}$. Every vertex of $C$ is within distance $3$ of $v_{i+3}$, and a vertex $w\in W$ is at distance at most $1+\min\{d_C(v_{i+3},x):x\in N(w)\cap C\}$ from $v_{i+3}$, which is at most $3$ unless the trace of $w$ is $\{v_i\}$. So there is $w_i\in C_i$ at distance $4$ from $v_{i+3}$. Every neighbour $z\in W$ of $w_i$ lies in $C_i$: if $z$ had a neighbour among $v_{i+2},v_{i+3},v_{i+4}$, then $d(w_i,v_{i+3})\le3$; if $z$ had the neighbour $v_{i-1}$ or $v_{i+1}$, then $w_i,v_i,v_{i\pm1},z$ would be a $4$-cycle. So $N(w_i)\subseteq\{v_i\}\cup C_i$.

\emph{Edges in $W$.} Let $x,y\in W$ be adjacent.
\begin{itemize}
\item $x\in C_0$, $y\in C_{\pm1}$: $x,v_0,v_{\pm1},y$ is a $4$-cycle.
\item $x\in C_0$, $y\in C_2$: $w_5,v_5,v_0,x,y,v_2,v_3,w_3$ is an induced $P_8$.
\item $x\in C_0$, $y\in C_3$: $w_1,v_1,v_0,x,y,v_3,v_4,w_4$ is an induced $P_8$. By symmetry, vertices of $C_i$ and $C_j$ with $i\ne j$ are never adjacent.
\item $x$ with trace $\{v_0,v_1\}$ and $y\in C_j$: for $j\in\{0,1\}$ the vertices $v_0,v_1,x,y$ span a $4$-cycle; for $j=5$ (resp.\ $j=2$) the cycle $x,v_0,v_5,y$ (resp.\ $x,v_1,v_2,y$) is a $4$-cycle; for $j=3$, $w_5,v_5,v_0,x,y,v_3,v_2,w_2$ is an induced $P_8$, and $j=4$ is symmetric.
\item $x,y$ with consecutive-pair traces $\{v_0,v_1\}$ and $\{v_j,v_{j+1}\}$: $j=0$ gives two common neighbours; $j=\pm1$ and $j=\pm2$ give $4$-cycles as above; $j=3$ gives the induced $P_8$ $w_2,v_2,v_1,x,y,v_4,v_5,w_5$.
\end{itemize}
In each displayed path the non-edges follow from the traces and from $N(w_i)\subseteq\{v_i\}\cup C_i$. So every edge inside $W$ joins two vertices of the same class $C_i$, and these edges form a matching (two neighbours in $C_i$ of one vertex would be two common neighbours with $v_i$).

\emph{Mass.} Let $s=\sum_i|C_i|$, let $t$ be the number of vertices with a consecutive-pair trace, and let $m$ be the number of edges inside $W$. Then $n=6+s+t$, $e(G)=6+s+2t+m$, and $T(G)=t+m$. So $S=2e-3T=12+2s+t-m=2n-t-m\le2n$.
\end{proof}


\begin{corollary}[proved, with a finite classification]\label{cor:r4}
REAL-133 holds for every connected strong $C_4$-free graph with $r(G)=4$ and $\lr(G)\le4$, and also for every one with $r(G)\ge5$ and $\lr(G)\le4$. Moreover $\RAD(4)$ holds.
\end{corollary}

\begin{proof}
Let $r=4$. By Theorem~\ref{thm:ess}, $\pp\ge7$. If $\pp=7$, Theorem~\ref{thm:r4} gives $\lr\le2$ and $\pp\ge r+2$. If $\pp\ge8$, then $\pp\ge r+4\ge r+\lceil\lr\rceil$. For $r\ge5$, $\pp\ge2r-1\ge r+4$.
\end{proof}

\subsection{Radius five, six and seven}
The next three theorems follow the pattern of Theorem~\ref{thm:r4}, iterated: Theorem~\ref{thm:cs} produces a dominating core, which itself has the radius-extremal structure one level lower, and the vertices outside it are controlled by trace classifications. We give outlines; the complete arguments, all trace lists and their independent recomputations are in the supplementary material.

\begin{theorem}[$\RAD(5)$; computer-assisted]\label{thm:rad5}
If $G$ is connected, strong $C_4$-free, of radius $5$ and $P_{10}$-free, then $e(G)=n+T(G)$, so $S(G)=2n-T(G)\le2n$ and $\lr\le2$.
\end{theorem}

\begin{proof}[Proof outline]
Theorem~\ref{thm:cs} with $k=10$ gives a connected dominating $H$ that is $C_{10}$ or $P_8$-free. In the latter case $r(H)=4$, $\pp(H)=7$, and the proof of Theorem~\ref{thm:r4} shows that $H=C_8$ or $H$ has a dominating $C_6$ with singleton or consecutive-pair traces and same-class matchings. (i) A dominating $C_{10}$: an outside vertex has neither a singleton nor a consecutive-pair trace (each gives a $P_{10}$), so its trace contains $c_0,c_d$ with $3\le d\le5$; the shortcut gives $\mathrm{ecc}(c_0)\le4$. So $G=C_{10}$. (ii) A dominating $C_8$: traces are singletons or consecutive pairs; witnesses at distance $5$ from the antipodal cycle vertex have singleton traces and are independent, which gives an induced $C_8$ with a pendant vertex at each cycle vertex. The classification has $24$ trace types meeting $C_8$, and adjacent outside vertices lie in a common class; counting gives $e=n+T$. (iii) A dominating $C_6$: witnesses give an induced $C_6$ with a pendant $P_2$ at each vertex; the classification has $31$ single and $31$ adjacent-pair trace types, every vertex with empty trace has exactly one parent, and edges between such vertices join vertices with the same parent; again $e=n+T$.
\end{proof}

The bound is attained by $C_8$ with a pendant vertex at each cycle vertex. The strong identity $e=n+T$ fails for every $r\ge6$ (Example~\ref{ex:radstar}), but the weaker bounds needed for $\RAD(r)$ survive.

\begin{theorem}[$\RAD(6)$; computer-assisted]\label{thm:rad6}
If $G$ is connected and strong $C_4$-free with $r(G)=6$ and $\pp(G)=11$, then $S(G)\le5n-36$.
\end{theorem}

\begin{proof}[Proof outline]
Theorem~\ref{thm:cs} with $k=12$ gives $H=C_{12}$ (then $G=C_{12}$ by the dominating-cycle shortcut argument) or a $P_{10}$-free $H$ with $r(H)=5$, $\pp(H)=9$, which has the structure from Theorem~\ref{thm:rad5}. Witnesses at distance $6$ extend this to an induced \emph{flower} $F_{L,t}$: a cycle $C_L$ with a pendant path on $t$ vertices at each cycle vertex, $(L,t)\in\{(6,3),(8,2),(10,1)\}$. Put $A=N(C)\setminus C$ and $R=V\setminus(C\cup A)$. Complete finite classifications give the single traces ($43$, $41$, $31$ types) and adjacent pairs ($55$, $41$, $11$), and show that there is no induced path $a\,x\,y\,z$ with $a\in A$ and $x,y,z\in R$, and that every vertex of $R$ has at most two neighbours in $A$. A written mass lemma bounds each component $J$ of $G[R]$: if every vertex has at most $d$ parents and no parent starts an induced $P_4$ into $J$, then $S(J)+2e(A,J)-3\tau\le(2d+1)|J|$, where $\tau=T_{A,J,J}$ is the number of triangles with one vertex in $A$ and two in $J$. With $S(G[C\cup A])\le2|C\cup A|$ and $d=2$ this gives $S\le5n-36$.
\end{proof}

\begin{theorem}[$\RAD(7)$; computer-assisted]\label{thm:rad7}
If $G$ is connected and strong $C_4$-free with $r(G)=7$ and $\pp(G)=13$, then $S(G)\le5n-42$, with equality if and only if $G=C_{14}$.
\end{theorem}

\begin{proof}[Proof outline]
The structure from Theorem~\ref{thm:rad6} propagates one level (the \emph{lifting lemma}) and gives a flower $F_{L,t}$ with $L+2t=14$, $(L,t)\in\{(12,1),(10,2),(8,3),(6,4)\}$, or the cycle $C_{14}$. The single traces ($37$, $51$, $57$, $55$ types) and adjacent pairs ($13$, $51$, $73$, $91$) are classified; rooted paths of the form $A\,R^k$ and vertices of $R$ with three parents in $A$ are excluded. For $(6,4)$ three further facts hold: a two-parent vertex has empty trace and nonadjacent parents with supports $c_i,c_{i+1}$; it starts no induced $P_4$ in $G[R]$; and two such vertices are never at distance $2$. Hence the two-parent vertices of each component of $G[R]$ form a clique of order at most $3$, so its mass is at most $5|J|$. Tip components save $L(t-1)$, giving $S\le5n-54$ for flowers, and the cycle gives $5n-42$.
\end{proof}

These structural statements are the first instances of a uniform structure theorem, $\mathrm{Struct}(r)$, stated precisely in Appendix~\ref{app:struct} (Definition~\ref{def:struct} and Theorem~\ref{thm:struct}): a connected strong $C_4$-free graph of radius $r\ge5$ with $\pp=2r-1$ is $C_{2r}$ or contains an \emph{admissible} induced flower $F_{2s,t}$ ($s+t=r$, $s\ge3$, $t\ge1$; admissibility includes the ambient distance inequalities $d_G(c_{i+s},z_i)\ge r$ for the arm tips $z_i$), and the vertices outside the flower satisfy explicit trace, matching, rooted-path and two-parent restrictions. The base case $\mathrm{Struct}(5)$ is supplied by the proof of Theorem~\ref{thm:rad5}; the lifting $\mathrm{Struct}(r)\Rightarrow\mathrm{Struct}(r+1)$ has a uniform written proof. The numerical bounds of Theorems~\ref{thm:rad5}--\ref{thm:rad7} alone would not suffice for this: the lifting uses the structural conclusions, in particular the distance inequalities, which cannot be imposed on arbitrary induced templates. What $\mathrm{Struct}(r)$ does not provide is control of the mass of the components of $G[R]$ hanging from the flower; that is the role of the rooted bounds of the next section.

\section{Rooted mass bounds and sharp \texorpdfstring{$\RAD(r)$}{RAD(r)} for \texorpdfstring{$r\le10$}{r<=10}}\label{sec:root}

Recall that $\pp_s(F)$ is the largest order of an induced path of $F$ starting at $s$, $D=d_F(s)$ and $m=e(F[N(s)])$.

\begin{theorem}[$\ROOT(q)$ for $q\le8$]\label{thm:root}
Let $F$ be connected and strong $C_4$-free, $s\in V(F)$ and $3\le q\le8$. If $\pp_s(F)\le q$, then
\[
S(F)\le(q-1)(|F|-1)-(q-3)D-m.
\]
For $q=3$ this is Lemma~\ref{lem:rootedp4} \textup{(}proved\textup{)}; for $q=4,5$ it is proved, with written proofs in the supplement \textup{(}\nolinkurl{lanes/astra_r3b/radK2/scratch/RAD6_ROOT_REPORT.md} and \nolinkurl{lanes/astra_r4/rad7/scratch/RAD7_ROOT5_REPORT.md}, the latter a charging argument with at most two charges per parent; the terms $-m$ are added by the amalgamation argument below\textup{)}; for $6\le q\le8$ it is computer-assisted. For each $q$, with the leading coefficient $q-1$ fixed, the coefficients $q-3$ and $1$ are optimal.
\end{theorem}

\begin{table}[ht]
\centering\small
\begin{tabular}{@{}cllp{7.2cm}@{}}
\toprule
$q$ & bound & status & main finite ingredients\\
\midrule
3 & $2(n-1)-m$ & proved & none\\
4 & $3(n-1)-D-m$ & proved & none\\
5 & $4(n-1)-2D-m$ & proved & component charging (cap $2$)\\
6 & $5(n-1)-3D-m$ & comp.-assisted & $P_6$: $17$ direct $+9$ bridge types, $406$ templates; $P_7$: $24+11$ types, $240$ obstruction certificates\\
7 & $6(n-1)-4D-m$ & comp.-assisted & near-parent shortcut; $332{,}231$ four-parent and $1{,}493{,}449$ six-parent completions\\
8 & $7(n-1)-5D-m$ & comp.-assisted & $\ge7$ parents; parents within distance $2$ of every $P_9$; $6{,}106$ direct-parent families; $8$ certificates\\
\bottomrule
\end{tabular}
\caption{$\ROOT(q)$ for $q\le8$ ($n=|F|$).}\label{tab:root}
\end{table}

\begin{proof}[Proof outline]
\emph{Component form.} A \emph{parent system} is a connected graph $J$ with $h=|J|$, a nonempty independent set $B$ of parents, each with a neighbour in $J$, and edges from $B$ to $J$ such that each vertex of $J$ has at most one parent and the whole graph is strong $C_4$-free. Its boundary mass is $M_B(J)=S(J)+2e(B,J)-3\tau$, where $\tau=T_{B,J,J}$ counts triangles with one vertex in $B$ and two in $J$. (Edges between parents enter neither $M_B(J)$ nor the hypothesis below, so they may be present or deleted.) Statement $C_q$ says: if, for every $u\in B$, every induced path of $J+u$ starting at $u$ has at most $q-1$ vertices, then $M_B(J)\le(q-1)h$. Closed-neighbourhood amalgamation of many copies of $F$ at the root (which preserves $\pp_s$ and $C_4$-freeness), together with auxiliary roots and summation over components, shows that $C_q$, $\ROOT(q)$ and the corrected form with the terms $-(q-3)D-m$ are equivalent. The same amalgamation applied to universal-root matching graphs proves optimality of the corrections.

\emph{Minimum counterexamples.} In a counterexample to $C_q$ with $h$ minimum, every nonempty $U\subseteq V(J)$ has deletion loss $\Delta_B(U)\ge(q-1)|U|+1$ (components of $J-U$ that lose all parents are bounded by $\ROOT(q-2)$ after a rooted prefix). Averaging the corrected $\ROOT(q-1)$ over the parents proves $C_q$ whenever there are at most $q-2$ parents, so a counterexample has at least $q-1$ parents. For $q=8$, each parent has at least two children, and $M_B(J)>7h$ and $e(B,J)\le h$ give $S(J)>5h$, so $\GEN_4$ (Theorem~\ref{thm:gen4}) forces an induced $P_9$ in $J$. The proof then excludes, by complete finite enumeration of the possible traces of additional parents, every way in which the parents can meet such a path. For $q=8$, all parents lie within distance $2$ of every induced $P_9$, and the configurations with $6,5,\dots,0$ direct parents on a $P_9$ are excluded in turn.
\end{proof}

\begin{theorem}[sharp $\RAD(r)$ for $5\le r\le10$; computer-assisted: uses Theorem~\ref{thm:root} and $\mathrm{Struct}(r)$]\label{thm:radsharp}
Let $5\le r\le10$ and let $G$ be connected and strong $C_4$-free with $r(G)=r$ and $\pp(G)=2r-1$. Then
\[
S(G)\le(r-1)n-2r(r-3),
\]
with equality only for $C_{2r}$. If $G$ is not a cycle, then $S(G)\le(r-1)n-\delta_r$ with $\delta_r=32,60,96,126,156,186$ for $r=5,\dots,10$.
\end{theorem}

\begin{proof}
By Theorem~\ref{thm:struct}, either $G=C_{2r}$, which has $n=2r$, $S=4r=(r-1)n-2r(r-3)$, or $G$ contains an admissible flower $F_{L,t}$ with $L=2s$, $s+t=r$, $s\ge3$, $t\ge1$, satisfying (S1)--(S5) of Definition~\ref{def:struct}. Let $C$, $A$, $R$ be as there and $n_0=|C\cup A|$. The first arm vertices lie in $A$, so $n_0\ge2L$, and the flower has $L(t+1)$ vertices, so $n\ge L(t+1)$.

\emph{Decomposition.} No vertex of $R$ is adjacent to $C$, and two $A$-neighbours of an $R$-vertex are nonadjacent by (S5). Hence every triangle meeting $R$ has one vertex in $A$ and two in $R$, and
\[
S(G)=S(G[C\cup A])+\sum_J M_A(J),
\]
the sum over the components $J$ of $G[R]$, with $M_A(J)$ the boundary mass of Section~\ref{sec:root} for the actual parents of $J$. By (S1)--(S2) the $A$-vertices have singleton or consecutive-pair traces and the edges of $G[A]$ join equal singleton traces and form a matching; counting edges and triangles gives $S(G[C\cup A])\le2n_0$.

\emph{Case $t=1$.} Then $R=\varnothing$ by (S4), $S\le2n_0=2n$, and $n\ge2L=4(r-1)$. So $(r-1)n-S\ge(r-3)n\ge4(r-1)(r-3)$.

\emph{Case $t\ge2$.} Put $\beta=t$. Since $t\le r-3\le7$, Theorem~\ref{thm:root} gives $\ROOT(t+1)$, which in component form is $M_B(J)\le t|J|$ whenever every parent $u\in B$ has $\pp_u(J+u)\le t$ and every vertex of $J$ has at most one parent in $B$. By (S3) every parent of $J$ satisfies the rooted condition; by (S5) the parents split into two classes $A_0,A_1$ such that each vertex of $J$ has at most one parent in each. Applying the component bound to the class that meets $J$ (either one if both do) and charging the other class at most $2$ per vertex of $J$ gives $M_A(J)\le(\beta+2)|J|$. For each $i$, the arm vertices $v_{i,2},\dots,v_{i,t}$ lie in one component $J_i$ of $G[R]$, these $L$ components are distinct, and all parents of $J_i$ have the singleton trace $\{c_i\}$ (this follows from the distance inequality for $z_i$; Appendix~\ref{app:struct}). So $J_i$ has unique parents and $M_A(J_i)\le\beta|J_i|$. Moreover the tip $z_i$ is a leaf or a degree-two triangle vertex of $G$: along a geodesic $x_0,\dots,x_{t-1}=z_i$ in $J_i+x_0$ from a parent $x_0$, an extra neighbour $w$ of $z_i$ must meet the geodesic (else the rooted path extends), a contact before $x_{t-3}$ shortens the distance, a contact at $x_{t-3}$ gives a $4$-cycle, so $w\sim x_{t-2}$, and two such $w$ give a $4$-cycle. Deleting $z_i$ therefore costs at most $2$, and $M_A(J_i)\le\beta(|J_i|-1)+2$. Summing, with $\sum_i|J_i|\ge L(t-1)$,
\[
S\le2n_0+(\beta+2)(n-n_0)-2L(t-1)-L(\beta-2)\le(\beta+2)n-L(3\beta+2t-4).
\]
With $\beta=t$ and $n\ge L(t+1)$ the deficit from $(r-1)n$ is at least $(s-3)n+2s(5t-4)\ge2s\bigl[(s+2)t+s-7\bigr]$.

\emph{Comparison.} For $5\le r\le10$ the minimum of these non-cycle deficits over $3\le s\le r-1$ is $\delta_r$ as stated (attained at $s=r-1$, $t=1$ for $r=5,6$, and at $s=3$ for $7\le r\le10$), and $\delta_r>2r(r-3)=20,36,56,80,108,140$. So every non-cycle satisfies the bound strictly and equality holds only for $C_{2r}$.
\end{proof}

For $r=8$ the strict form of $\ROOT(6)$ improves the non-cycle constant to $132$ (supplement, \nolinkurl{lanes/astra_r7/radpush}). For $r=5,6,7$ the inequality (but not the equality statement) also follows from Theorems~\ref{thm:rad5}--\ref{thm:rad7} and $n\ge2r$. For $r\ge11$ the same computation needs $\ROOT(q)$ for $q\le r-2$, and gives the conditional bound below.

\begin{corollary}[consequence of Theorems~\ref{thm:A}, \ref{thm:B}, \ref{thm:gen3}, \ref{thm:gen4}, \ref{thm:r4}, \ref{thm:rad5}; the second sentence also uses Theorem~\ref{thm:radsharp}]\label{cor:five}
REAL-133 holds for every connected graph with $\lr(G)\le5$. Moreover $\RAD(r)$ holds for $4\le r\le10$, so alternative~\textup{(b)} of Theorem~\ref{thm:B} cannot occur when $K\le9$.
\end{corollary}

\begin{proof}
Theorem~\ref{thm:A} covers $\lr\le3$. For $3<\lr\le5$, $K=\lceil\lr\rceil-1\in\{3,4\}$. Alternative~(a) of Theorem~\ref{thm:B} is excluded by $\GEN_3$ and $\GEN_4$. Alternative~(b) has $r=K+1\in\{4,5\}$ and $\pp=2r-1$, and Theorems~\ref{thm:r4} and~\ref{thm:rad5} give $\lr\le2<r-1$. The last statement follows from Corollary~\ref{cor:r4} and Theorem~\ref{thm:radsharp}, since $(r-1)n-2r(r-3)<(r-1)n$.
\end{proof}

\paragraph{Beyond $q=8$.} The following partial results are verified. Here an \emph{outer component} is a component of $F-N[s]$, and its parents are its neighbours in $N(s)$.
\begin{itemize}
\item \emph{Corrected $\ROOT(9)$ with few parents} (computer-assisted, via $\ROOT(8)$): if $\pp_s(F)\le9$ and every outer component has at most seven parents, then $S(F)\le8(n-1)-6D-m$.
\item \emph{Ordinary $\ROOT(9)$ for small order} (computer-assisted): if $\pp_s(F)\le9$ and $|F|\le190$, then $S(F)\le8(|F|-1)$. This is the weaker inequality without the terms $-6D-m$; the amalgamation argument that adds those terms uses arbitrarily many copies of $F$ and therefore does not apply to a statement restricted to bounded order. The relaxation used in the proof first permits an excess at order $191$; no graph violating $\ROOT(9)$ is known.
\item \emph{Finite local criterion} (proved): for $q\ge4$, $\ROOT(q-1)$ together with a finite decidable statement $\mathcal L_q(b)$ (Definition~\ref{def:Lq} in Appendix~\ref{app:Lq}) implies $\ROOT(q)$. For $q=9$ all $350$ seeds of $\mathcal L_9$ are open under the depth-zero rules.
\item \emph{Conditional uniform $\RAD$} (proved): if $\ROOT(q)$ holds for $3\le q\le r-2$, then every connected strong $C_4$-free graph with radius $r\ge5$ and $\pp=2r-1$ satisfies $S\le(r-1)n-\min\{2r(r-3),\,30r-114\}$. This is the computation in the proof of Theorem~\ref{thm:radsharp}: for $t\ge2$ the minimum of $2s[(s+2)t+s-7]$ over $3\le s\le r-2$ is attained at an endpoint, with values $30r-114$ and $6(r-2)(r-3)$, and the case $t=1$ gives $4(r-1)(r-3)$.
\end{itemize}

\section{Restricted classes: all indices}\label{sec:classes}

\subsection{Blocks that are bipartite or cycles}
Let $\mathcal B$ be the class of connected strong $C_4$-free graphs in which every block is bipartite or a cycle (edges are blocks; cycles may be odd). It contains every connected strong $C_4$-free bipartite graph and every connected strong $C_4$-free cactus, but of course not graphs containing a $4$-cycle: the restriction matters, since $K_{2,3}$ rooted in its two-vertex part has $\pp_s=3$ and $S=12>2(n-1)-m=8$. A rooted graph $(G,s)$ with $\pp_s(G)=q$ and $\mathrm{ecc}(s)=q-1$ is called \emph{terminal}; for it put $R=n-q+1-D\ge0$, the number of vertices outside $N[s]$ and outside a root-to-farthest-vertex geodesic. Set $c_2=0$, $c_3=2$, $c_4=3$ and $c_q=2\lfloor(q-1)/2\rfloor$ for $q\ge5$.

\begin{theorem}[\textup{(1)} and \textup{(2)} proved; \textup{(3)} computer-assisted, as it uses $\mathrm{Struct}(r)$]\label{thm:blocks}
\begin{enumerate}
\item Every terminal instance in $\mathcal B$ satisfies $S(G)\le2(D+q-2)-m+c_qR$.
\item Every $F\in\mathcal B$ with $\pp_s(F)\le q$ satisfies $S(F)\le c_q(n-1)-(c_q-2)D-m$; in particular $\ROOT(q)$ holds in $\mathcal B$ for every $q\ge3$. If every block is bipartite or a triangle, one may take $c_q=2\lfloor(q-1)/2\rfloor$ for all $q\ge3$.
\item \textup{(}uses $\mathrm{Struct}(r)$\textup{)} For every $r\ge5$, every $G\in\mathcal B$ with $r(G)=r$ and $\pp(G)=2r-1$ satisfies $S\le(r-1)n-2r(r-3)$, with equality exactly for $C_{2r}$; if $G$ is not a cycle, $S\le(r-1)n-4(r-1)(r-3)$, with equality exactly for $C_{2r-2}$ with one pendant vertex at each cycle vertex.
\end{enumerate}
\end{theorem}

\begin{proof}
\emph{Two terminal facts.} An endpoint $x$ of a longest rooted induced path is not a cut vertex: a neighbour of $x$ in another component of $G-x$ would extend the path. If $z$ is at distance $q-1$ from $s$ and $v_0=s,\dots,v_{q-1}=z$ is a geodesic, any neighbour $y\ne v_{q-2}$ of $z$ must meet an earlier path vertex (else the rooted path extends); meeting $v_i$ with $i\le q-4$ shortens $d(s,z)$, and meeting $v_{q-3}$ gives a $C_4$; so $y\sim v_{q-2}$, and two such $y$ would give a $C_4$. Hence $L(z)=2d(z)-3t(z)\le2$. Moreover, if $x\ne s,z$, $G-x$ is connected and $z$ survives, then $q-1\le d_{G-x}(s,z)\le\pp_s(G-x)-1\le q-1$, so the instance stays terminal with the same $q$.

\emph{Endpoint losses.} Let $x$ be the endpoint of a longest rooted path of order $q$. It lies in a single block. If that block is a triangle the loss $L(x)$ is $1$, if an edge $2$, if a cycle of length at least $5$ it is $4$. If it is bipartite, let $\ell$ be the order of the final path segment inside it. Block-tree separation shows that the neighbours of $x$ in this block can only be blocked by earlier path vertices in the same block; every neighbour other than the predecessor has a blocker in the colour class of $x$, the vertex two positions earlier cannot block (a $C_4$), and each position blocks at most one neighbour. So $d(x)\le\lfloor(\ell-1)/2\rfloor\le\lfloor(q-1)/2\rfloor$, there are no triangles at $x$, and $L(x)\le c_q$ for $q\ge5$.

\emph{Bases.} For $q\le4$ a nontrivial bipartite block would yield, by the same endpoint count, a rooted induced path of order at least $5$, so these instances are cacti. For a $C_4$-free cactus with $c_3$ triangle blocks and $c_L$ longer cycle blocks, $S=2(n-1)+2c_L-c_3$ and $m\le c_3$. For $q=3$ no longer cycle occurs, and (1) follows. For $q=4$ every longer cycle is a $C_5$ through $s$. A chosen length-$3$ geodesic uses at most one edge of each such $C_5$: if it used two edges and then left the cycle, replacing that two-edge arc by the other three-edge arc would give an induced rooted path of order $5$. So each $C_5$ contributes two vertices to $R$, $R\ge2c_L$, and $[2(D+2)-m+3R]-S=R-2c_L+c_3-m\ge0$. For $q=2$ the root is universal and $S=2D-m$.

\emph{Induction on $n$.} Let $E$ be the set of endpoints of rooted induced paths of order $q$ and $z$ a farthest vertex, so $z\in E$. If some $x\in E\setminus\{z\}$ exists, delete it: $x\notin N[s]$, $x$ is not a cut vertex, its loss is at most $c_q$, the instance remains terminal with the same $q$, $D$ and $m$, and $R$ drops by one. If $E=\{z\}$, delete $z$: the new rooted maximum is $q-1$, the eccentricity $q-2$, $R$ is unchanged and the loss is at most $2$; induction gives $S\le2(D+q-3)-m+c_{q-1}R+2\le2(D+q-2)-m+c_qR$. Deletion preserves membership in $\mathcal B$. This proves (1); the bipartite-or-triangle variant is identical.

(2) Attach a new path with $q-1$ edges at $s$. The result is terminal of index $q$, lies in $\mathcal B$, has order $n+q-1$, mass $S+2(q-1)$ and root degree $D+1$, and $m$ is unchanged. Substituting into (1) gives (2); since $c_q\le q-1$, $\ROOT(q)$ follows.

(3) By Theorem~\ref{thm:struct}, $G=C_{2r}$ or $G$ contains an admissible flower $F_{L,t}$, $L=2s$, $s+t=r$, with (S1)--(S5); we use the decomposition and notation of the proof of Theorem~\ref{thm:radsharp}. In $\mathcal B$ the block containing the even cycle $C$ is $C$ itself or a bipartite block; a vertex with a consecutive-pair trace would put a triangle into that block, so there are no consecutive-pair traces and all parents have singleton traces, and the auxiliary rooted graphs used to bound the components hanging from the flower stay in $\mathcal B$. Part (2) with leading coefficient $t$ bounds one parent parity class, the other costs at most $2|J|$, and (1) applied to the selected tip components $J_i$ gives a saving of $(t-1)(t-2)$ each. Summing, $S\le(t+2)n-Lt(t+1)$, and with $n\ge L(t+1)$ the deficit from $(r-1)n$ is at least $(s-3)n+2st(t+1)\ge2s(r-s+1)(r-3)$. For $t\ge2$, that is $3\le s\le r-2$, the function $2s(r-s+1)$ is concave in $s$ with endpoint values $6(r-2)>4(r-1)$, so the deficit exceeds $4(r-1)(r-3)$. For $t=1$, $R=\varnothing$ by (S4), $S\le2n$ and $n\ge2L=4(r-1)$, so the deficit is at least $(r-3)n\ge4(r-1)(r-3)$. Equality forces $n=4(r-1)$, so $V(G)=V(F_{L,1})$; the arm vertices have distinct singleton traces, so by (S2) there are no further edges and $G=F_{2r-2,1}$, the cycle $C_{2r-2}$ with one pendant vertex at each cycle vertex, which indeed has $S=2n$. Finally $4(r-1)(r-3)>2r(r-3)$ for $r\ge5$, and $C_{2r}$ attains $2r(r-3)$.
\end{proof}

\begin{lemma}[endpoint blockers; proved]\label{lem:blockers}
Let $G$ be strong $C_4$-free and triangle-free, and let $x_0x_1\cdots x_k$ be an induced path that cannot be extended at $x_k$ \textup{(}for instance a longest induced path\textup{)}. If $d(x_k)\ge2$, then $k\ge d(x_k)+1$. If moreover $G$ is bipartite, then $k\ge2d(x_k)$.
\end{lemma}

\begin{proof}
Let $y\ne x_{k-1}$ be a neighbour of $x_k$. It is not on the path (the path is induced) and, since the path cannot be extended, $y$ has a neighbour $x_i$ with $i\le k-1$. Here $i\ne k-1$ (a triangle) and $i\ne k-2$ (the $4$-cycle $x_{k-2}x_{k-1}x_ky$), so $i\le k-3$. A vertex $x_i$ is adjacent to at most one neighbour of $x_k$ (Lemma~\ref{lem:basic}(1)). Hence the $d(x_k)-1\ge1$ such $y$ need distinct blockers among $x_0,\dots,x_{k-3}$, and $k-2\ge d(x_k)-1$. If $G$ is bipartite, $y$ has the colour opposite to $x_k$, so its blocker $x_i$ has the colour of $x_k$, that is $i\equiv k\pmod 2$ and $i\le k-4$; there are $\lfloor k/2\rfloor-1$ such positions, so $\lfloor k/2\rfloor\ge d(x_k)$.
\end{proof}

\begin{proposition}[proved]\label{prop:regbip}
Every connected regular bipartite strong $C_4$-free graph satisfies $\GEN_K$ for every $K$ and REAL-133. Every Moore graph \textup{(}a $d$-regular graph of diameter $2$ and girth $5$\textup{)} satisfies REAL-133.
\end{proposition}

\begin{proof}
Let $G$ be $d$-regular. If $d=1$ then $G=K_2$, with $r=1$, $\lr=1$, $\pp=2$ and $S=n$, and both statements hold. Let $d\ge2$. As $G$ is triangle-free, $\lambda\equiv d$ and $\lr=d$. Lemma~\ref{lem:blockers} applied to a longest induced path gives $\pp\ge2d+1$, and Theorem~\ref{thm:ess} gives $\pp\ge2r-1$; so $\pp\ge r+d$ both when $r\le d+1$ and when $r\ge d+1$. If $\pp\le2K$, then $2d+1\le2K$, so $S=dn<Kn$; thus $\GEN_K$ holds. A Moore graph of degree $d\ge2$ is triangle-free and strong $C_4$-free with $r=2$ and $\lr=d$, and the first part of Lemma~\ref{lem:blockers} gives $\pp\ge d+2=r+\lr$.
\end{proof}

\begin{remark}\label{rem:moore}
A Moore graph of degree $57$, if it exists, has $n=57^2+1=3250$ and $S=57n$. It satisfies REAL-133 by Proposition~\ref{prop:regbip}, which only needs $\pp\ge59$. The statement $\GEN_{57}$ imposes nothing on it (its conclusion $S\le57n$ holds with equality), but $\GEN_{56}$ does: since $S>56n$, $\GEN_{56}$ requires $\pp\ge2\cdot56+1=113$. We do not know whether such a graph would have $\pp\le112$. If it did, it would refute $\GEN_{56}$ while satisfying REAL-133; so it \emph{could} separate $\GEN_K$ from REAL-133, but its existence alone would refute neither.
\end{remark}

\subsection{Bipartite graphs and \texorpdfstring{$\GEN_5$}{GEN5}}

\begin{theorem}[computer-assisted]\label{thm:bip}
\begin{enumerate}
\item No nonempty bipartite strong $C_4$-free $P_{11}$-free graph with colour classes $X,Y$ has $d(x)\ge4$ for all $x\in X$ and $d(y)\ge3$ for all $y\in Y$.
\item Every bipartite strong $C_4$-free $P_{11}$-free graph with $n\ge2$ and colour classes of sizes $a,b$ satisfies $S=2e\le5n-8-|a-b|$. In particular $\GEN_5$ holds for bipartite graphs.
\end{enumerate}
\end{theorem}

\begin{proof}[Proof outline]
(1) The degree assumptions force an induced $P_9$ starting in $X$ (a written argument with a two-record finite check). An extension certificate then starts from this $P_9$: each record is a labelled induced graph with a vertex below its required degree, and lists a child for \emph{every} admissible trace of a new neighbour (all subsets of the appropriate colour class containing that vertex, subject only to strong $C_4$-freeness and $P_{11}$-freeness). The certificate has $1030$ records ($436$ terminal), $1029$ arcs and largest order $27$; every arc raises the order by one. (2) Property (1) is hereditary. Repeatedly delete an $X$-vertex of degree at most $3$ or a $Y$-vertex of degree at most $2$ (taking $X$ to be the smaller class), charging $3$ or $2$ edges; the last two vertices span at most one edge, which leaves enough unused allowance for the bound.
\end{proof}

The bound extends to clique-sums: if $G$ is strong $C_4$-free, $P_{11}$-free and is assembled from connected bipartite pieces of order at least $2$ and $t$ triangles by $s_1$ identifications along a vertex and some identifications along an edge, with imbalances $d_i$ of the bipartite pieces, then $S(G)\le5n-8-4t-3s_1-\sum_id_i$ (proved from Theorem~\ref{thm:bip}). The $P_{12}$-free analogue of Theorem~\ref{thm:bip}(2) fails (the incidence graph of the affine plane of order~$3$).

\section{The band \texorpdfstring{$5<\lr\le6$}{5<l<=6}}\label{sec:cells}

\begin{proposition}[consequence]\label{prop:cells}
A counterexample to REAL-133 with $5<\lr\le6$ has $(r,\pp)=(4,9)$ \textup{(}Cell~A\textup{)} or $(r,\pp)=(5,10)$ \textup{(}Cell~B\textup{)}.
\end{proposition}

\begin{proof}
Here $K=5$. Alternative (b) of Theorem~\ref{thm:B} ($r=6$, $\pp=11$) is excluded by Theorem~\ref{thm:rad6}. In (a), a violation needs $\pp\le r+5$ and $\pp\ge2r-1$. For $r\le3$, $\pp\le8$ and $\GEN_4$ gives $\lr\le4$. For $r=4$, $\pp\in\{7,8,9\}$; $\pp=7$ is Theorem~\ref{thm:r4} and $\pp=8$ is $\GEN_4$. For $r=5$, $\pp\in\{9,10\}$, and $\pp=9$ is Theorem~\ref{thm:rad5}.
\end{proof}

Both cells would follow from $\GEN_5$, or from $\GEN_5$ restricted to these radii, and both are closed for bipartite graphs by Theorem~\ref{thm:bip}. Neither cell is closed in general. Appendix~\ref{app:cells} gives the precise setting (the nested cores produced by Theorem~\ref{thm:cs}, the subcases A1, A2, B1, parents and component masses) and the partial closures that were verified, each with its exact hypotheses.

\section{Negative results}\label{sec:neg}

All examples below are explicit graphs. They were checked by direct computation, and the arguments are given.

\begin{example}[the identity $e=n+T$ fails for $r\ge6$]\label{ex:radstar}
Let $s\ge3$, $t\ge3$ and $m\ge1$ be integers. Let $G_{s,t,m}$ be $C_{2s}=c_0\cdots c_{2s-1}$ with a pendant path of $t$ edges at every cycle vertex, together with $m$ internally disjoint paths $c_0xzyc_1$ (\emph{ears}). Then $n=2s(t+1)+3m$, $e=2s(t+1)+4m$ and $T=0$, so $e-n-T=m\ge1$; for $(s,t,m)=(3,3,1)$, $n=27$, $e=28$, $r=6$, $\pp=11$. These graphs satisfy REAL-133 ($\lr=2e/n<8/3$), but they show that Theorem~\ref{thm:rad5} has no direct analogue for $r\ge6$, because $r(G_{s,t,m})=s+t$ and $\pp(G_{s,t,m})=2(s+t)-1$.

\emph{Proof.} Two vertices have at most one common neighbour, so $G$ is strong $C_4$-free, and it is triangle-free. An ear has length $4>d(c_0,c_1)$, so it shortens no distance between old vertices; hence the eccentricities of the flower vertices are unchanged, and $c_0$, which is within distance $2$ of every ear vertex, still has eccentricity $s+t$. An ear vertex $w$ reaches the tip $z$ of the arm at $c_s$ only through $c_0$ or $c_1$, so $d(w,z)\ge1+\min\{d(c_0,z),d(c_1,z)\}=s+t$. So $r=s+t$. For the path order, let $B$ be the block formed by $C$ and the ears. An induced path $P$ meets $B$ in a subpath $Q$, and can leave $B$ only at an end of $Q$ that is a cycle vertex, through at most $t$ arm vertices there. Split according to $Q\cap\{c_0,c_1\}$. If $Q$ avoids both, $Q$ lies in $c_2\cdots c_{2s-1}$ or inside one ear, so $|Q|\le2s-2$. If $Q$ contains exactly one of them, say $c_0$, then $Q-c_0$ consists of at most two pieces, each inside one ear ($\le3$ vertices, ending at an ear vertex) or inside $c_{2s-1}\cdots c_2$; so either $Q\subseteq C$ with $|Q|\le2s-1$, or $|Q|\le3+1+(2s-2)$ with at most one cycle end, or $|Q|\le7$ with no cycle end. If $Q$ contains both, they are consecutive on $Q$, each side continues into an ear ($\le2$ vertices) or into the cycle, and the same three possibilities occur with $|Q|\le2s-1$, $|Q|\le2s+2$ and $|Q|\le6$. Hence $|P|\le\max\{2s-1+2t,\ 2s+2+t,\ 7\}=2s+2t-1$, using $t\ge3$ and $s\ge3$; $C$ minus one vertex together with the two arms at its ends attains this.
\end{example}

\begin{example}[star bounds fail for $K=5$]\label{ex:star5}
The natural analogue of step (iii) in the proof of Theorem~\ref{thm:gen4} for $K=5$ (``$L\ge6$ and $P_{11}$-free exclude an induced $K_{1,6}$'') is false. Two Petersen graphs sharing a vertex give $n=19$, $S=60$, a triangle-free strong $C_4$-free graph with $\min L=6$, $\pp=9$ and $\lambda=6$ at the shared vertex. For $b\ge2$, gluing $b$ Petersen graphs along a common edge gives a $2$-connected graph $B_b$ with $n=8b+2$, $S=28b+2$, $\min L=6$, $\pp=9$ and $\lambda=2b+1$ at the two hubs ($B_1$ is the Petersen graph, with $\pp=5$). So singleton losses cannot bound $\lambda$. These graphs violate the set-loss condition of a minimum counterexample (the eight non-hub vertices of one Petersen block have loss $28<5\cdot8+1$), so they are not counterexamples to $\GEN_5$.
\end{example}

\begin{example}[endpoint inequalities]\label{ex:ep}
Let $G$ be connected and strong $C_4$-free and let $(u,v)$ be an oriented edge. Consider the induced paths whose first two vertices are $u,v$ in this order; let $q$ be their maximum order and $E$ the set of last vertices of those of order $q$. The \emph{endpoint inequality} $\mathrm{EP}(q)$ asserts $D_G(E)\le q|E|$ for all such $(G,(u,v))$ with maximum $q$. For a root $s$ with $\pp_s(G)=q$ let $E$ be the set of last vertices of the induced paths of order $q$ starting at $s$; the rooted variant $\mathrm{EP}^+(q)$ asserts $D_G(E)\le(q-1)|E|$. The equivalence $\mathrm{EP}(q)\Leftrightarrow\mathrm{EP}^+(q-1)$ for $q\ge4$ is proved, and $\mathrm{EP}^+(6)$ holds (computer-assisted: a discharging argument with a $265$-record local certificate), which gives $\mathrm{EP}(7)$.

$\mathrm{EP}(9)$ fails. Let $G$ have vertices $0,\dots,14$ and adjacency lists
\[
\begin{array}{llll}
0{:}\ 4,8,11 & 1{:}\ 5,8,10 & 2{:}\ 6,8,9 & 3{:}\ 7,10\\
4{:}\ 0,9,10 & 5{:}\ 1,9,11 & 6{:}\ 2,10,11 & 7{:}\ 3,11\\
8{:}\ 0,1,2,12 & 9{:}\ 2,4,5 & 10{:}\ 1,3,4,6,13 & 11{:}\ 0,5,6,7,14\\
12{:}\ 8,13,14 & 13{:}\ 10,12 & 14{:}\ 11,12. &
\end{array}
\]
It is triangle-free and strong $C_4$-free, with $n=15$, $e=23$, $S=46$, graph6 code \verb|N?`@?b_JAsOw?G?A_?o|. For the oriented edge $(12,8)$ the maximum order is $q=9$, attained only by $12,8,2,9,4,10,3,7,11$ and $12,8,2,9,5,11,7,3,10$; so $E=\{10,11\}$ and $D_G(E)=20>18$. $\mathrm{EP}^+(8)$ fails on the subgraph induced by $0,\dots,11$ (graph6 \verb|K?`@?b_JAsOw|) with root $8$: $\pp_8=8$, $E=\{10,11\}$ and $D(E)=16>14$. Componentwise charging versions fail as well. $\mathrm{EP}(8)$ is open.
\end{example}

\begin{example}[the Petersen graph refutes a dominating-core budget]\label{ex:cb}
Let $k=\pp-r$. Call an induced subgraph $H$ of $G$ an \emph{eligible core} if it is proper, connected, dominating, and $\pp(H)\le\pp(G)-2$ (Theorem~\ref{thm:cs} supplies such cores, apart from a cycle alternative). The budget $\mathrm{CB}$ asks for an eligible core with $S(G)-S(H)\le k(n-|H|)+|H|$; together with REAL-133 for $H$ by induction it would imply REAL-133. The Petersen graph has $S=30$, $\pp=5$, $r=2$, $k=3$. An eligible core has $\pp(H)\le3$; a connected triangle-free strong $C_4$-free $P_4$-free graph is a star; a star in the Petersen graph with $k\le3$ leaves dominates at most $4+2k$ vertices, which is $10$ only for $k=3$, so the eligible cores are exactly the ten closed neighbourhoods $N[v]$. Each has $S(H)=6$ and $S(G)-S(H)=24>3\cdot6+4=22$. Among all $28{,}003{,}330$ eligible cores of connected strong $C_4$-free graphs with $n\le12$, only these ten fail. A repaired version restricted to $\lr>5$ remains open (Conjecture~\ref{conj:cbhigh}).
\end{example}

\begin{example}[a scalar relaxation that does not imply the conjecture]\label{ex:heawood}
The degree, triangle and distance data of the Heawood graph ($n=14$, cubic, $T=0$, $r=3$, $S=42$, $Z=28$) satisfy all the identities used in our counting arguments, including the Moore-type identity $\sum_vd(v)^2=n(n-1)+6T-2Z$ ($Z$ the number of nonadjacent pairs without a common neighbour), together with the generic lower bounds $\pp\ge2r-1$ and $\pp\ge\min_v\lambda(v)+2$, if $\pp$ is replaced by the value $5$. But $5<r+\lr=6$. So this particular system of scalar constraints together with these two path lower bounds does not imply REAL-133. This does not exclude arguments that use further inequalities relating degrees, eccentricities and triangles to $\pp$. (The real Heawood graph has $\pp=7$.)
\end{example}

\begin{example}[radius-extremal graphs with large $\lr$]\label{ex:plane}
Let $F$ be the incidence graph of the projective plane $\mathrm{PG}(2,q)$, with $f=|F|=2(q^2+q+1)$ and $S(F)=(q+1)f$, rooted at a vertex. Let $t=f-1$ and $b\ge1$. Take $C_6=c_0\cdots c_5$, attach a pendant path with $t$ edges at each cycle vertex, and attach $b$ copies of $F$ at $c_0$ by identifying their roots with $c_0$. The result is strong $C_4$-free (each new block meets the rest in one cut vertex) and
\[
n=6f+b(f-1),\qquad S=12f+b(q+1)f,\qquad r=t+3,\qquad \pp=2t+5=2r-1.
\]
Indeed the distances inside the flower are unchanged; every vertex of an attached copy is within distance $f-1\le t$ of $c_0$, so each cycle vertex has eccentricity $t+3$ and the flower already has radius $t+3$, while a vertex $x$ of a copy has distance $d_F(x,c_0)+t+3$ from the tip of the arm at $c_3$. An induced path meets the cycle in at most five consecutive vertices, and each of its two ends outside the cycle lies in one arm or one copy of $F$, with at most $t$ vertices; the flower contains an induced path on $2t+5$ vertices. Mass is additive at cut vertices. For fixed $q$, $\lr\to(q+1)f/(f-1)$ as $b\to\infty$; letting then $q\to\infty$, with $r=f+2\sim2q^2$, this limit is $\sim\sqrt{r/2}$. For example $q=5$, $b=17$ gives $n=1409$, $r=64$, $\pp=127$ and $S=7068>5n$. These graphs satisfy $\RAD(r)$, but they show that $\lr$ is unbounded among radius-extremal graphs, so no absolute bound such as $\lr\le5$ holds for all $r$, and $\RAD(r)$ needs a bound depending on $r$.
\end{example}

Further refuted routes (path-count recurrences, random induced walks, ball-deletion lemmas, and two approaches to TD) are recorded in Appendix~\ref{app:routes}.

\section{Conjectures and open problems}\label{sec:open}

By Theorem~\ref{thm:B} and Corollary~\ref{cor:five}, REAL-133 follows from $\GEN_K$ for all $K\ge5$ (or only in the form needed in alternative (a)) together with $\RAD(r)$ for all $r\ge11$. The following are open.

\begin{enumerate}
\item \textbf{The two cells of Proposition~\ref{prop:cells}}, equivalently REAL-133 for $5<\lr\le6$. Open subcases are listed in Appendix~\ref{app:cells}.
\item \textbf{$\GEN_K$ for $K\ge5$.} On the endpoint route, the missing steps are $\mathrm{EP}(8)$ and a boundary-selection inequality for $K=5,6$ (inequality~(8) of \nolinkurl{lanes/astra_r5/ep8/scratch/REPORT.md}); $\mathrm{EP}(9)$ already fails (Example~\ref{ex:ep}), so this route cannot be uniform in $q$. Scalar component linear programs are provably insufficient.
\item \textbf{$\ROOT(q)$ for $q\ge9$}, and hence $\RAD(r)$ for $r\ge11$. For $q=9$, all $350$ seeds of $\mathcal L_9$ remain open.
\item \textbf{Universal $\mathrm{EB}_6$ and unrestricted triangle-free $\GEN_6$.} After R35--R40 the avoiding class A38 is not emptied; high-vertex mass bounds (H31/H32/H35/H36) and the Robertson puncture reduction are accepted, but the residual mass bound remains open. Triangle-inclusive $\GEN_6$/$T_6$, the finite bound F42, general $\GEN_K$, $\RAD(r)$ for $r\ge11$, and REAL-133 itself remain open.
\end{enumerate}

\begin{conjecture}[uniform rooted and radius bounds]\label{conj:uniform}
$\ROOT(q)$ holds for every $q\ge3$, and for every $r\ge5$ the bound $S\le(r-1)n-2r(r-3)$ of Theorem~\ref{thm:radsharp} holds, with equality only for $C_{2r}$.
\end{conjecture}

The conjecture is proved for $q\le8$ and $r\le10$ (Theorems~\ref{thm:root} and~\ref{thm:radsharp}) and for all $q$ and $r$ in the class $\mathcal B$ (Theorem~\ref{thm:blocks}). It would follow from the next statement.

\begin{conjecture}[terminal correction TD]\label{conj:td}
If $F$ is connected and strong $C_4$-free, $\pp_s(F)=q$ and $\mathrm{ecc}(s)=q-1$, then
\[
S(F)\le(q-1)(n-1)-(q-3)D-m-(q-2)(q-3).
\]
\end{conjecture}

Equivalently, $2\beta-(T-m)\le(q-3)\,|V\setminus(P\cup N[s])|$ for a root-to-farthest geodesic $P$, where $\beta=e-T-n+1$. It is proved that TD implies $\ROOT(q)$ for all $q$ and Conjecture~\ref{conj:uniform}. TD holds in $\mathcal B$ (Theorem~\ref{thm:blocks}(1)) and has no violation among the $3{,}646{,}151$ terminal roots of connected strong $C_4$-free graphs with $n\le14$.

\begin{conjecture}[$\mathrm{CB}_{\mathrm{high}}$]\label{conj:cbhigh}
Every connected strong $C_4$-free graph $G$ that is not a cycle and has $\lr(G)>5$ has an eligible core $H$ with $S(G)-S(H)\le(\pp-r)(n-|H|)+|H|$.
\end{conjecture}

It is proved that $\mathrm{CB}_{\mathrm{high}}$, together with Corollary~\ref{cor:five} and an order induction, implies REAL-133 without any $\GEN_K$. There is no test data: no connected strong $C_4$-free graph with $n\le12$ has $\lr>5$. Example~\ref{ex:cb} shows that the restriction is necessary, and Example~\ref{ex:heawood} shows that a proof cannot rely on the scalar constraints listed there alone.

\section{Method: computer assistance and verification}\label{sec:method}

The results of Sections~\ref{sec:gen3}--\ref{sec:neg} and of the appendices were obtained in rounds. In each round the AI model GPT~6~Astra worked on a precisely stated target in a fresh session, with a context file listing the definitions and the previously accepted results, and wrote a report together with its certificates as JSON files. A claim was accepted only after a second, separate GPT~6~Astra session, acting as referee, had checked it. In this paper ``referee'' and ``independent'' refer to such separate sessions of the same model; they are not external peer review and not an independent human assessment. The referee session shared no session state with the prover beyond the claim, the prover's report and the certificate files. It was instructed to be adversarial, to recompute the finite certificates with code it wrote itself (without running the prover's code as a substitute), and to return a verdict \textsc{correct}, \textsc{correct with fixes}, \textsc{incorrect} or \textsc{unverified} for each claim. Only claims judged \textsc{correct} or \textsc{correct with fixes} appear here as theorems, with the fixes incorporated. Separate sessions of one model can share systematic failure modes; this is why we emphasise inspectable written proofs, independently written executable checks, and, for $\GEN_3$, kernel verification, rather than the session protocol itself. For the early rounds (Sections~\ref{sec:gen3}--\ref{sec:rad}), the referee audits were complemented by stand-alone checkers using only the Python standard library and by separately written C and Python re-implementations of the finite checks; for example, the finite parts of Lemmas~\ref{lem:forks}--\ref{lem:mixed} were reproduced by three implementations, and the $\GEN_4$ certificate by a C verifier with its own trace enumeration.

The referees did reject or correct claims. Examples: a lemma in the cactus case of TD was false as stated (a $7$-vertex counterexample) and was repaired; a reported count of $4{,}881$ maximal parent families was in fact $2{,}218$ (search leaves had been counted); a rooted-$P_5$ statement proved only for $n\le12$ was demoted to a conjecture until a later round proved it; several statements needed explicit hypotheses (ambient strong $C_4$-freeness, actual parents meeting the component, $t\ge2$) without which they are false, and these hypotheses are included above. Search totals from runs that could not be replayed, and statements about the implementation history that the referees could not authenticate (marked \textsc{unverified} in the verdict files), are not used in any proof and are not presented here as independently checked; where an original script was missing, the referee's replacement computation is the one supplied and cited.

The certificate types are: extension DAGs (Theorems~\ref{thm:gen4} and~\ref{thm:bip}), in which every record lists a child for every admissible trace of a forced new neighbour, and every arc raises the order by one, so that the walk inside a hypothetical counterexample terminates; complete trace and adjacent-pair classifications on fixed cores (Sections~\ref{sec:rad} and~\ref{sec:root}); and exclusion lists with explicit $C_4$ and induced-path witnesses. In all of them, completeness of the enumeration of traces is the critical point, and it is what the referees recomputed. Apart from $\GEN_3$, none of the computer-assisted results is formally verified.

A later \emph{continuous proving window} (24~October~2026, ending at 18:00 Asia/Shanghai, concurrent Codex cap~$2$) ran structural rounds R35--R40 on $\mathrm{EB}_6$ and triangle-free $\GEN_6$. Each accepted claim again required a fresh Astra referee session returning \textsc{correct} or \textsc{correct with fixes}; round certificates \nolinkurl{ROUND_BRANCH_R35}--\texttt{R40} record verdict SHA-256 hashes and claim breakdowns. Only those accepted claims appear in Theorem~G and Appendix~\ref{app:window}; UNVERIFIED open targets and provenance narratives are excluded.


\section{Data availability and reproducibility}\label{sec:data}
The directory \texttt{data/} accompanying this manuscript contains (see \nolinkurl{data/README.md}):
\begin{itemize}
\item \nolinkurl{check_lemmas.py}, a stand-alone standard-library Python checker for the finite classifications of Lemma~\ref{lem:c6}, Fact~2 in Lemma~\ref{lem:p6} (also proved by hand) and Lemma~\ref{lem:c8}, for Proposition~\ref{prop:theta}, and for small-order checks;
\item \texttt{rounds/}: the round certificates \nolinkurl{ROUND_BRANCH_R1}--\texttt{R40} (claims, statuses and SHA-256 hashes of the payloads), the stand-alone checkers \nolinkurl{verify_round_branch_r1}--\texttt{r6.py}, the $\GEN_4$ extension certificate, the referee verdicts and the repairs written by the referees, and the project notes \nolinkurl{PROOF_NOTES.md};
\item \nolinkurl{rounds/lanes/}: for every prover and referee session of rounds 1--10, the complete report (\texttt{report.md} and, where present, \nolinkurl{scratch/REPORT.md}), the task statement, and all scripts and certificate payloads produced, including the $\ROOT(6)$--$\ROOT(8)$ exclusion certificates, the bipartite extension certificate (\nolinkurl{lanes/astra_r8/cells/scratch/bipartite_p11_certificate.json}), the fixed-core certificates of the cells and the endpoint certificates; files larger than $5$\,MB are compressed with \texttt{xz} (the SHA-256 hashes of the uncompressed files are in \nolinkurl{lanes/COMPRESSED_FILES.sha256}); nothing was omitted, and the bundle has about $280$\,MB;
\item \texttt{lean/}, the Lean~4 development of Section~\ref{sec:lean} with its build and axiom logs;
\item \texttt{certificate/} and \texttt{src/}, consistency data from earlier versions, no longer used in proofs.
\end{itemize}
Appendix~\ref{app:index} lists, for each result, the written proof, the certificates and the referee report. The checkers \nolinkurl{verify_round_branch_r*.py} locate their inputs relative to their own directory and write their reports to the directory given by the environment variable \texttt{OUTDIR} (default: the working directory), so they run from a clean extraction, for example
\begin{center}
\nolinkurl{mkdir out && OUTDIR=out python3 data/rounds/verify_round_branch_r2.py}
\end{center}
They use only the Python~3 standard library; rounds 1 and 2 take a few minutes, the others seconds. Some of them recompute (rounds 1 and 2 recompute the $\GEN_3$ configuration counts and re-verify the full $\GEN_4$ certificate); others only check saved metadata (for example the round-6 checker tests that a saved \texttt{violations} field is zero, which does not rerun that census). \nolinkurl{data/README.md} says which is which, and gives the commands for the referees' replacement computations, their resource requirements and the pinned dependencies (Python~3 standard library, gcc, nauty~2.8.9, Lean v4.35.0 with the Mathlib revision fixed in \nolinkurl{lake-manifest.json}). A permanent archive identifier will be assigned when the supplement is deposited.

\section*{Acknowledgments}
The author used AI models and AI tools in this work, as follows. The AI model GPT~6~Astra produced, in separate prover sessions, the proofs of the local lemmas of Section~\ref{sec:gen3}, of $\GEN_4$ and of the results of Sections~\ref{sec:rad}--\ref{sec:neg} and the appendices, together with their finite certificates and the counterexamples; in earlier work it also suggested the hand proof of Fact~2 in Lemma~\ref{lem:p6}, the equality case of Lemma~\ref{lem:p6}, and the argument of Theorem~\ref{thm:r4} with the classification in Lemma~\ref{lem:c8}. Separate GPT~6~Astra sessions acted as referees: they checked every accepted result, recomputed the certificates with their own code, and supplied several repairs; GPT~6~Astra sessions also reviewed drafts of this manuscript. The AI model GPT~6.1~Sol, run through OpenAI Codex, was used as a simulated referee of earlier versions of the manuscript. Grok Bot, an AI agent tool, was used to orchestrate the sessions, to run computations and the stand-alone checkers, to produce the Lean formalisation of $\GEN_3$, and to draft and edit the text. OpenAI Codex was used as the tool through which the GPT~6~Astra and GPT~6.1~Sol sessions were run. The author takes full responsibility for the content of this paper, including any errors.

\appendix

\section{The structure theorem \texorpdfstring{$\mathrm{Struct}(r)$}{Struct(r)}}\label{app:struct}

Let $L=2s\ge6$ and $t\ge1$. The \emph{flower} $F_{L,t}$ consists of a cycle $C=c_0c_1\cdots c_{L-1}c_0$ and, for every $i$, an arm $c_i=v_{i,0},v_{i,1},\dots,v_{i,t}=z_i$; indices of cycle vertices are taken modulo $L$, and $\delta(i,j)$ denotes the distance of $c_i,c_j$ on $C$. For a vertex $x\notin V(F)$ its \emph{trace} is $N(x)\cap V(F)$. Write $S_{i,j}=\{v_{i,j}\}$, $E_{i,j}=\{v_{i,j},v_{i,j+1}\}$ and $B_i=\{c_i,c_{i+1}\}$.

\begin{definition}\label{def:struct}
An induced copy of $F_{L,t}$ in a graph $G$ is \emph{admissible} if $G$ is strong $C_4$-free and $P_{L+2t}$-free and
\begin{equation}\label{eq:D}
d_G(c_{i+s},z_i)\ge s+t\qquad(0\le i<L).
\end{equation}
For an admissible flower put $A=N(C)\setminus V(C)$ and $R=V(G)\setminus(V(C)\cup A)$; the \emph{parents} of a component $J$ of $G[R]$ are its neighbours in $A$, and $J+a=G[V(J)\cup\{a\}]$. For $r\ge5$, $\mathrm{Struct}(r)$ is the statement: \emph{every connected strong $C_4$-free graph $G$ with $r(G)=r$ and $\pp(G)=2r-1$ is $C_{2r}$ or contains an admissible $F_{2s,t}$ with $s+t=r$, $s\ge3$, $t\ge1$, such that}
\begin{itemize}
\item[(S1)] every vertex outside the flower has trace $\varnothing$, $S_{i,j}$ or $E_{i,j}$ $(0\le j<t)$, or $B_i$; in particular every vertex of $A$ has a singleton or consecutive-pair trace on $C$;
\item[(S2)] every edge of $G[A]$ joins two vertices with the same singleton trace on $C$, and these edges form a matching;
\item[(S3)] for every component $J$ of $G[R]$ and every parent $a$ of $J$, every induced path of $J+a$ starting at $a$ has at most $t$ vertices;
\item[(S4)] if $t=1$, then $R=\varnothing$;
\item[(S5)] every vertex of $R$ has at most two neighbours in $A$; if it has two, it has empty trace, and its two parents are nonadjacent with traces $S_{i,0},S_{i+1,0}$ for some $i$ \textup{(}this occurs only when $t\ge3$\textup{)}. Consequently $A=A_0\cup A_1$, where $A_\epsilon$ contains the vertices with singleton trace $c_i$, $i\equiv\epsilon\pmod2$, and the vertices with a pair trace are put into either class, and every vertex of $R$ has at most one neighbour in each class.
\end{itemize}
\end{definition}

The classification of adjacent pairs of outside vertices is also part of the record: the admissible unordered trace pairs of two adjacent outside vertices are exactly $(\varnothing,\varnothing)$, $(S_{i,j},S_{i,j})$ for $j<t$, $(\varnothing,S_{i,j})$ and $(\varnothing,E_{i,j})$ for $j\le t-2$, $(\varnothing,B_i)$ for $t\ge2$, $(S_{i,j},S_{i,j+3})$ and $(E_{i,j},S_{i,j+3})$ for $j\le t-4$, and $(S_{i,j},E_{i,j+2})$ for $j\le t-3$; there are $1+L(2t+1)$ single traces (including $\varnothing$) and $1+L$, $1+5L$, $1+L(6t-9)$ adjacent pairs for $t=1$, $t=2$, $t\ge3$.

\begin{theorem}[base case computer-assisted; the rest proved in the supplement]\label{thm:struct}
\begin{enumerate}
\item If $G$ contains an admissible $F_{L,t}$, then \textup{(S1)--(S5)} and the pair classification hold.
\item $\mathrm{Struct}(5)$ holds.
\item $\mathrm{Struct}(r)\Rightarrow\mathrm{Struct}(r+1)$ for every $r\ge5$. Hence $\mathrm{Struct}(r)$ holds for every $r\ge5$.
\end{enumerate}
\end{theorem}

\begin{proof}[Sources and outline]
(1) and (3) are proved in \nolinkurl{lanes/astra_r5/radr/report.md}, \S\S1--7, and were audited in \nolinkurl{audits_r2_r6/astra_r6_auditStruct.report.md}. The single-trace classification rests on the observation that, for a vertex adjacent to $v_{i,a}$ and $v_{j,b}$, \eqref{eq:D} is preserved exactly when $|a-b|\le2$ ($i=j$) or $\delta(i,j)+|a-b|\le2$ ($i\ne j$), and that $C_4$-safety is also a pairwise condition; every excluded trace yields an explicit induced $P_{L+2t}$ obtained by joining two arms through a long arc of $C$. An arbitrary-length bridge lemma (an induced path outside $F$ with empty internal traces cannot join two different arms when one end has positive depth) gives (S3): an induced path $a,x_1,\dots,x_t$ with $a\in A$ and $x_j\in R$ could be completed by a cycle segment of $L-1$ vertices and a full arm to an induced $P_{L+2t}$. (S4) is the case $t=1$ of (S3), and (S5) follows from the pair classification and the bridge lemma. All lists were in addition recomputed for the $44$ flowers with $5\le s+t\le12$.

(2) The proof of Theorem~\ref{thm:rad5} gives, for a connected strong $C_4$-free graph of radius $5$ without $P_{10}$, that $G=C_{10}$, or $G$ contains $C_8$ with a pendant vertex $w_i$ at each $c_i$, or $C_6$ with a pendant $P_2$ at each $c_i$, where the $w_i$ (respectively the ends of the pendant paths) were chosen at distance $5$ from $c_{i+s}$. These are admissible copies of $F_{8,1}$ and $F_{6,2}$, and (1) supplies (S1)--(S5).

(3) Let $G$ have radius $r+1$ and no $P_{2r+2}$. Theorem~\ref{thm:cs} gives a connected dominating induced $H$ that is $C_{2r+2}$ (then $G=C_{2r+2}$ by the dominating-cycle shortcut argument) or $P_{2r}$-free; then $r(H)=r$ and $\pp(H)=2r-1$. If $H=C_{2r}$, witnesses at distance $r+1$ from the cycle vertices have singleton traces and give arms of depth $1$. Otherwise $\mathrm{Struct}(r)$ gives a flower in $H$, and a witness $w_i$ with $d_G(c_{i+s},w_i)\ge r+1$ lies outside $H$; for $t\ge2$ each of its $H$-neighbours lies in a component $J\subseteq R(H)$ all of whose parents have the singleton trace $c_i$ (for $t=1$ use (S1)--(S2) instead), by the same distance computation as below, and a shortest parent-to-neighbour path extended by $w_i$ is an arm of depth $t+1$. Edges between new tips are excluded by \eqref{eq:D} or by an explicit induced $P_{L+2t+2}$. So $G$ contains an admissible $F_{L,t+1}$, and (1) applies.
\end{proof}

\emph{Selected tip components.} The distance inequality~\eqref{eq:D} is what makes the flower usable for mass bounds. Let $t\ge2$, let $J_i$ be the component of $G[R]$ containing $v_{i,2},\dots,v_{i,t}$, and let $a$ be a parent of $J_i$ with $c_j\in N(a)$. By (S3) a shortest path from $a$ to $z_i$ in $J_i+a$ has at most $t-1$ edges, so $d_G(c_{i+s},z_i)\le\delta(i+s,j)+1+(t-1)\le s+t$, with equality only if $\delta(i+s,j)=s$, that is $j=i$. By~\eqref{eq:D}, every parent of $J_i$ has trace exactly $\{c_i\}$. In particular the $L$ components $J_i$ are distinct and have unique parents. These are the facts used in the proofs of Theorems~\ref{thm:radsharp} and~\ref{thm:blocks}(3). Radius is not hereditary, and cannot be imposed on the finite templates used in the classifications; inequality~\eqref{eq:D} is the hereditary substitute that the finite pruning uses.

\section{The finite local criterion \texorpdfstring{$\mathcal L_q(b)$}{Lq(b)}}\label{app:Lq}

We use the parent systems $(J,B)$, the boundary mass $M_B(J)=S(J)+2e(B,J)-3T_{B,J,J}$ and the statement $C_q$ of Section~\ref{sec:root}; $C_q$ is equivalent to $\ROOT(q)$. Assume $\ROOT(q-1)$ and let $(J,B)$ be a counterexample to $C_q$ with $h=|J|$ minimum. For $\varnothing\ne U\subseteq V(J)$ let $\Delta_B(U)=M_B(J)-M_B(J-U)$, where $M_B(J-U)$ is summed over the components of $J-U$ with their actual parents.

\begin{lemma}[proved; \nolinkurl{lanes/astra_r8/root8/scratch/REPORT.md}, \S\S2--3]\label{lem:Lqbasic}
\begin{enumerate}
\item $\Delta_B(U)\ge(q-1)|U|+1$ for every nonempty $U\subseteq V(J)$; in particular $L_J(v)=2d_J(v)-3t_J(v)\ge q-2$ for every $v$, and $(J,B)$ has at least $q-1$ parents.
\item \emph{Templates.} Let $K\subseteq V(J)$, $B_0\subseteq B$ nonempty, and let $Q$ be $G[K\cup B_0]$ together with an auxiliary root $\rho$ adjacent exactly to $B_0$. A set $A\subseteq V(Q)\setminus\{\rho\}$ is an \emph{admissible trace} if adding a new vertex $y$ with $N(y)=A$ gives a strong $C_4$-free graph with no induced path on $q+1$ vertices starting at $\rho$; let $\mathcal A(Q)$ be the set of admissible traces meeting $K$. For $\varnothing\ne U\subseteq K$ let $\delta_Q(U)=2\#\{\text{edges of }Q-\rho\text{ meeting }U\}-3\#\{\text{triangles of }Q-\rho\text{ meeting }U\}$, let $u_0(U)$ be the number of vertices of $U$ without a neighbour in $B_0$, and $c_U(A)=2|A\cap U|-3\#\{vw\in E(Q-\rho):v,w\in A,\{v,w\}\cap U\ne\varnothing\}$. If every $A\in\mathcal A(Q)$ with $c_U(A)>0$ has $|A|\ge2$, then
\[
\Delta_B(U)\le\delta_Q(U)+2u_0(U)+\sum_{A\in\mathcal A(Q)}\max\{0,c_U(A)\}.
\]
\item \emph{Raising rule.} If $Q$ is embedded as an induced subgraph in $(J,B)$ and $v\in K$ satisfies $L_K(v)-3T_{B_0,K,K}(v)<q-2$, then $v$ has a neighbour $y\in V(J)\setminus K$ whose trace on $Q$ is admissible, contains $v$, and misses $N_K(v)\cup\{\rho\}$.
\end{enumerate}
\end{lemma}

\begin{definition}\label{def:Lq}
The \emph{seeds} $\mathcal S_q$ are the templates formed by $q-1$ parents, one private neighbour (child) of each, and $\rho$, where the children induce a strong $C_4$-free graph on $q-1$ vertices other than $P_{q-1}$ (one representative per isomorphism class). A template $Q$ \emph{wins by deletion} if some nonempty $U\subseteq K$ satisfies the hypothesis of Lemma~\ref{lem:Lqbasic}(2) with right-hand side at most $(q-1)|U|$. For an integer $b\ge0$ define $W_q(Q,b)$ recursively: it holds if $Q$ wins by deletion, or if $b>0$ and there is a vertex $v$ as in Lemma~\ref{lem:Lqbasic}(3) such that $W_q(Q+y_A,b-1)$ holds for every admissible trace $A$ that contains $v$ and misses $N_K(v)\cup\{\rho\}$ (vacuously if there is none); a vertex $v$ as in~(3) with no such trace also gives a win at every depth. Otherwise $W_q(Q,b)$ is false. Finally $\mathcal L_q(b)$ is the statement that $W_q(Q,b)$ holds for every $Q\in\mathcal S_q$.
\end{definition}

$\mathcal L_q(b)$ is a finite decidable statement: retained templates have at most $2q-1+b$ vertices.

\begin{theorem}[proved; same source, \S5; referee report \nolinkurl{verdicts/astra_r8v/v_root8.report.md}]\label{thm:Lq}
For $q\ge4$ and $b\ge0$, $\ROOT(q-1)$ and $\mathcal L_q(b)$ together imply $\ROOT(q)$.
\end{theorem}

\begin{proof}
A minimum counterexample to $C_q$ has at least $q-1$ parents by Lemma~\ref{lem:Lqbasic}(1), hence contains a seed. Following the raising rule extends the induced embedding along the certificate; a deletion win contradicts Lemma~\ref{lem:Lqbasic}(1), and a raising vertex with no admissible trace contradicts Lemma~\ref{lem:Lqbasic}(3).
\end{proof}

The numbers of seeds are $7,17,43,116,350$ for $q=5,\dots,9$. No depths $b_q$ completing $\mathcal L_q(b_q)$ are known for $q\ge9$, and the criterion is only sufficient.

\section{The band \texorpdfstring{$5<\lr\le6$}{5<l<=6}: setting and partial closures}\label{app:cells}

Throughout, $G$ is connected and strong $C_4$-free, in Cell~A ($r=4$, $\pp=9$) or Cell~B ($r=5$, $\pp=10$). To \emph{close} a subcase means to prove $S(G)\le5n$ for every graph of the cell satisfying its hypotheses. All results below hold at every order. Sources are given in Appendix~\ref{app:index}.

\subsection*{Nested cores and subcases}
Theorem~\ref{thm:cs} gives a connected dominating induced subgraph $H$ (the \emph{core}) that is $C_{\pp+1}$ or $P_{\pp-1}$-free, and applying it again inside $H$ gives an \emph{inner core} $J$ dominating $H$; then $J$ $2$-dominates $G$.
\begin{itemize}
\item \textbf{A1} (Cell~A): $H$ is $P_8$-free with $r(H)=3$, and $J\subseteq H$ is a dominating $P_6$-free induced subgraph with $r(J)=2$.
\item \textbf{A2} (Cell~A): $G$ has an induced $C_6$ that $2$-dominates $G$ (for instance when $r(H)=4$, $\pp(H)=7$ and $H$ has a dominating induced $C_6$ as in the proof of Theorem~\ref{thm:r4}).
\item \textbf{B1} (Cell~B): $(r(H),\pp(H))=(4,8)$, $(r(J),\pp(J))=(3,6)$, and $G$ has no induced $2$-dominating $C_6$.
\end{itemize}
\emph{Layers for A2.} For a $2$-dominating induced $C=c_0\cdots c_5$ put $A=N(C)\setminus C$, $R=V(G)\setminus(C\cup A)$ and $L=G[C\cup A]$. For a component $D$ of $G[R]$ let $P=N(D)\cap A$ be its \emph{actual parents}, $h=|D|$, $k=|P|$, and
\[
M(D)=S(D)+2e(P,D)-3\tau_1-3\tau_2,
\]
where $\tau_1$ counts triangles with one vertex in $P$ and two in $D$, and $\tau_2$ triangles with two vertices in $P$ and one in $D$. Then $S(G)=S(L)+\sum_DM(D)$. A parent is \emph{normal} if its trace on $C$ is a singleton or a consecutive pair, and \emph{antipodal} if it is $\{c_i,c_{i+3}\}$; there are at most three antipodal vertices.

\emph{Layers for A1 and B1.} Put $L_0=V(J)$, $L_1=\{v:d(v,J)=1\}$, $L_2=\{v:d(v,J)=2\}$ (these are all vertices) and $H^*=G[L_0\cup L_1]$, the \emph{first-layer core}. A component $D$ of $G[L_2]$ has actual parents $N(D)\cap L_1$. \emph{Unique parents relative to $H^*$} means $|N(v)\cap V(H^*)|=1$ for every $v\in L_2$. This does not follow from the corresponding condition relative to $H$, and $\pp(H)\le8$ does not imply $\pp(H^*)\le8$ (counterexamples of order $2m+16$ and $22$ in the supplement); results stated under unique parents assume the condition relative to $H^*$ explicitly, together with any path bound on $H^*$ that they use.

\subsection*{Verified partial closures}
\begin{enumerate}
\item \emph{Bipartite graphs and clique-sums.} Both cells are closed for bipartite graphs (Theorem~\ref{thm:bip}) and for strong $C_4$-free $P_{11}$-free graphs assembled from connected bipartite pieces and triangles by vertex and edge identifications (the bound after Theorem~\ref{thm:bip}).
\item \emph{Dominating cycles.} In Cell~A, if $G$ has a dominating induced $C_k$ with $k\in\{6,7,8,10\}$, then $S\le5n$; in Cell~B the same holds for a dominating induced $C_9$ or $C_{11}$, and when $\pp(H)=7$ (computer-assisted).
\item \emph{$P_6$-free blocks.} Every $2$-connected strong $C_4$-free $P_6$-free graph has at most $13$ vertices, and there are exactly $11$ of them (proved, with a finite classification).
\item \emph{A2, few parents} (proved). If $h\ge2$ and $k\le4$, then $M(D)\le5h$; if $k\le3$, then $M(D)\le4h$. This follows from the identity
\[
(k+1)h-M(D)=\sum_{v\in D}\bigl[k-w_v+(b_v-1)(d_D(v)-1)\bigr]+\tau_1+3T(D)+3\tau_2,
\]
where $b_v=|N(v)\cap P|$ and $w_v=|(N(v)\cap P)\cup\bigcup_{z\in N_D(v)}(N(z)\cap P)|\le k$, in which every summand is nonnegative for $h\ge2$. The bound fails for singletons, where $M(\{v\})=2k-3e(G[P])$.
\item \emph{A2, five normal parents} (proved, with a $63$-trace finite classification on $P_6$). If $D$ has at least five normal parents, then $D$ is a connected graph whose blocks are edges and triangles, of diameter at most $4$; such a $D$ has a vertex of eccentricity at most $2$ in $D$.
\item \emph{A2, global charge} (proved). Select a set $\mathcal C$ of residual components each of which is a singleton or an edge/triangle block tree of diameter at most $4$, and assume $M(D)\le5|D|$ for every unselected component. Let $I$ be the bipartite incidence graph between $\mathcal C$ and the set $Q$ of their actual parents. If $I$ is a forest and every edge of $G[A]$ has a common neighbour on $C$, then $S(G)<5n$. More generally $S(G)\le5n-\Delta+2|Q|+2\beta(I)-2f(I)-\eta$, where $\Delta=5|L|-S(L)$, $\beta(I)$ and $f(I)$ are the cycle rank and the number of components of $I$, and $\eta\ge0$ is an explicit sum of component savings.
\item \emph{A2, independent first layer} (proved). If $A$ is independent and every residual component has at most three parents, then $S\le5n-12$.
\item \emph{Rooted-$P_5$ theorem} (proved). If $Q$ is $2$-connected, strong $C_4$-free and $P_7$-free and some vertex starts no induced $P_5$, then $r(Q)\le2$. So every vertex of such a graph of radius $3$ starts an induced $P_5$.
\item \emph{B1 inner cores with cut attachments} (proved). Let $J$ be a $2$-connected strong $C_4$-free $P_7$-free graph of radius $3$, let $G\supseteq J$ be strong $C_4$-free and $P_{11}$-free, and suppose that every component $D$ of $G-J$ is attached to $J$ at a single vertex $v_D$. Then $S(G)\le5n-2|J|-3\beta-\tau$, where $\beta=\sum_D|N(v_D)\cap D|$ and $\tau=\sum_De(G[N(v_D)\cap D])$.
\item \emph{Minimum counterexamples to $\GEN_5$ and $\GEN_6$} (computer-assisted). A counterexample to $\GEN_5$ of minimum order is $3$-connected, for every $3$-vertex cut $X$ the graph $G-X$ has a component consisting of one triangle-free vertex of degree $3$, and every vertex starts an induced $P_9$. A counterexample to $\GEN_6$ of minimum order is $4$-connected, with the analogous property for $4$-vertex cuts, and every vertex starts an induced $P_8$. A counterexample to $\GEN_K$ with radius $r\ge2$ has $n\ge K^2+r$.
\end{enumerate}
Further closures for special attachment patterns in A1 and B1 (prefix closures with unique outer parents, hub-pair dominated blocks, maximal-block reductions) are proved in the supplement under the hypotheses stated there; we do not restate them. \emph{Examples.} Families with $\lr\to4$ exist in both cells, and among graphs with $n\le15$ the largest values of $\lr$ are $2.93$ (Cell~A) and $2.67$ (Cell~B). The componentwise bound $M(D)\le5|D|$ fails in A2: $C_6$ with a leaf $a_i$ at each $c_i$, plus an edge $xy$ with $x\sim a_0,a_1$ and $y\sim a_2,a_3,a_4$, has $n=14$, $r=4$, $\pp=9$, $S=36$, and $D=\{x,y\}$ has $M(D)=12>10$; there is also an infinite three-parent family.

\emph{Open subcases.} A1 with non-unique or direct attachments; A2 components with five to seven parents of which at most four are normal, and cyclic incidence structures not meeting the budget of item~6; B1 when the maximal block containing the hub core has $\pp\in\{9,10\}$ and every one of its vertices starts an induced $P_7$.

\section{Other refuted routes}\label{app:routes}
Here $I_j$ is the number of \emph{oriented} induced paths on $j$ vertices (each path counted once in each direction). The path-count recurrence $I_{2k+1}\ge(\lr-k)I_{2k-1}$ holds for $k=1$ in the form $I_3\ge2(\lr-1)n$, with equality for cycles (for unoriented counts it fails already for $C_5$), but fails for $k=2$: the Petersen graph with a handle and $29$ leaves ($n=42$) has $I_3=940$, $I_5=220$ and $nI_5-(S-2n)I_3=-160$. For the random induced walk, the start vertex is chosen uniformly and each step chooses uniformly among the neighbours of the current end that keep the path induced, until there is none; the proposed inequality $\mathbb E\lceil|P|/2\rceil\ge\lr$ is tight on the Petersen graph but fails on the Robertson graph, where $\mathbb E\lceil|P|/2\rceil=71/18<4=\lr$. Single-ball and two-ball deletion lemmas fail on the Heawood graph and on $B_5$ (Example~\ref{ex:star5}), respectively. For the terminal correction TD, singleton deletion and geodesic amplification provably cannot work: Petersen amalgams with a tail attain TD with equality at $q=5$ with the tip the unique cheap non-cut vertex, and $C_4$-safe geodesic amalgamation can increase the rooted maximum.

\section{Index of proofs, certificates and referee reports}\label{app:index}
Paths are relative to \texttt{data/rounds/}; \texttt{L/} abbreviates \texttt{lanes/}, and \texttt{V/} abbreviates \texttt{verdicts/}. For each lane, the written proof is in \texttt{report.md} and, where present, \nolinkurl{scratch/REPORT.md} (or the named file); the certificates and the scripts that produced them are in its \texttt{scratch/} directory, listed in its \nolinkurl{MANIFEST.json} where present. Referee lanes (\nolinkurl{L/astra_r*v/v_*}) contain the referees' own code and replacement computations.

\begingroup\small
\begin{longtable}{@{}>{\raggedright\arraybackslash}p{3.6cm}>{\raggedright\arraybackslash}p{5.8cm}>{\raggedright\arraybackslash}p{4.6cm}@{}}
\toprule
Result & Proof and certificates & Referee / audit\\
\midrule
\endhead
Lemmas \ref{lem:forks}--\ref{lem:mixed}, Thm.~\ref{thm:gen3} & \S\ref{sec:gen3}; \nolinkurl{L/astra_r1/A}; \nolinkurl{verify_round_branch_r1.py}; \texttt{data/lean/} & \nolinkurl{L/astra_r2/audit}; \nolinkurl{audits_r2_r6/astra_r2_audit.report.md}\\
Thm.~\ref{thm:gen4} & \nolinkurl{L/astra_r2/gen4}; \nolinkurl{referee_data/GEN4_star_certificate_p9.json}; \nolinkurl{verify_round_branch_r2.py} & \nolinkurl{audits_r2_r6/astra_r3_audit4.report.md}\\
Lemma~\ref{lem:c8}, Thm.~\ref{thm:r4} & \S\ref{sec:r4}; \nolinkurl{data/check_lemmas.py} & --\\
Thm.~\ref{thm:rad5} & \nolinkurl{L/astra_r3/rad5/scratch/RAD5_PROOF.md} & \nolinkurl{audits_r2_r6/astra_r4_audit5.report.md}\\
Thm.~\ref{thm:rad6}, $\ROOT(4)$ & \nolinkurl{L/astra_r3b/radK2/scratch/RAD6_ROOT_REPORT.md} & \nolinkurl{audits_r2_r6/astra_r4_audit6.report.md}\\
Thm.~\ref{thm:rad7}, $\ROOT(5)$ & \nolinkurl{L/astra_r4/rad7/scratch/RAD7_ROOT5_REPORT.md}; \nolinkurl{referee_data/r7/ROOT5_EQUALITY_PROOF_radpush.md} & \nolinkurl{audits_r2_r6/astra_r5_auditR7.report.md}; \nolinkurl{V/astra_r7v/v_radpush}\\
Thm.~\ref{thm:struct} & \nolinkurl{L/astra_r5/radr} & \nolinkurl{audits_r2_r6/astra_r6_auditStruct.report.md}\\
$\ROOT(6)$, sharp $\RAD(8)$ & \nolinkurl{L/astra_r7/root6}, \nolinkurl{L/astra_r7/radpush}; \nolinkurl{referee_data/r7/P6_EQUALITY_REPAIR.md} & \nolinkurl{V/astra_r7v/v_root6}, \texttt{v\_radpush}\\
$\ROOT(7)$, sharp $\RAD(9)$ & \nolinkurl{L/astra_r8/root7} (certificates \nolinkurl{near_P7.json}, \nolinkurl{near_P8.json}, \nolinkurl{near_P9.json}, \nolinkurl{prefix_*.json}, \nolinkurl{*_exclusions*.json}, \dots) & \nolinkurl{V/astra_r8v/v_root7}\\
$\ROOT(8)$, sharp $\RAD(10)$, $\ROOT(9)$ partial & \nolinkurl{L/astra_r9/root8b} (eight proof certificates \nolinkurl{far3_parentcap14.json}, \nolinkurl{near_dK_a7*.json}, $K=0,\dots,6$) & \nolinkurl{V/astra_r9v/v_root8b}\\
Thm.~\ref{thm:Lq}, Lemma~\ref{lem:Lqbasic} & \nolinkurl{L/astra_r8/root8} (\S\S2--5; \nolinkurl{far4_patches.json}, \nolinkurl{child_seeds8.json}) & \nolinkurl{V/astra_r8v/v_root8}\\
Thm.~\ref{thm:radsharp} equality; component equivalence; conditional $\RAD$ & \S\ref{sec:root}; \nolinkurl{L/astra_r9/unify} \S\S1,3 & \nolinkurl{V/astra_r9v/v_unify}\\
Thm.~\ref{thm:blocks}, TD in $\mathcal B$ & \S\ref{sec:classes}; \nolinkurl{L/astra_r10/td}; \nolinkurl{referee_data/r9/cactus_proof_counterexample.json} & \nolinkurl{V/astra_r10v/v_td}\\
Thm.~\ref{thm:bip}, clique-sums & \nolinkurl{L/astra_r8/cells} (\nolinkurl{bipartite_p11_certificate.json}, \nolinkurl{bipartite_proof.py}); \nolinkurl{L/astra_r9/cellsA2} & \nolinkurl{V/astra_r8v/v_cells}, \nolinkurl{V/astra_r9v/v_cellsA2}\\
App.~\ref{app:cells}, items 2, 3 & \nolinkurl{L/astra_r4/cells}, \nolinkurl{L/astra_r5/cellsre}, \nolinkurl{L/astra_r7/cellsA1A2B1} (fixed-core certificates \nolinkurl{A_*_stars.json}, \nolinkurl{A_*_finish.json}, \dots) & \nolinkurl{audits_r2_r6/astra_r5_cellsre.report.md}; \nolinkurl{V/astra_r7v/v_cellsA1A2B1}\\
App.~\ref{app:cells}, items 4--7 & \nolinkurl{L/astra_r8/cells}, \nolinkurl{L/astra_r8/cellsA1B1}, \nolinkurl{L/astra_r9/cellsA2} & \nolinkurl{V/astra_r8v/v_cells}, \texttt{v\_cellsA1B1}; \nolinkurl{V/astra_r9v/v_cellsA2}\\
App.~\ref{app:cells}, items 8, 9 & \nolinkurl{L/astra_r9/cellsA1B1b} & \nolinkurl{V/astra_r9v/v_cellsA1B1b}\\
App.~\ref{app:cells}, item 10 & \nolinkurl{L/astra_r8/sel8}, \nolinkurl{L/astra_r8/ep8}, \nolinkurl{L/astra_r9/sel9} (endpoint certificates \nolinkurl{critical*.json}, \nolinkurl{star_K5.json}); \nolinkurl{L/astra_r10/cb} & \nolinkurl{V/astra_r8v/v_sel8}, \texttt{v\_ep8}; \nolinkurl{V/astra_r9v/v_sel9}; \nolinkurl{V/astra_r10v/v_cb}\\
Ex.~\ref{ex:radstar}, \ref{ex:plane} & \S\ref{sec:neg}; \nolinkurl{L/astra_r5/radr} \S10; \nolinkurl{verify_round_branch_r3.py} & \nolinkurl{audits_r2_r6/astra_r6_auditStruct.report.md}\\
Ex.~\ref{ex:star5} & \nolinkurl{L/astra_r3/genK} & \nolinkurl{verify_round_branch_r3.py}\\
Ex.~\ref{ex:ep}, $\mathrm{EP}^+(6)$ & \nolinkurl{L/astra_r4/ep}, \nolinkurl{L/astra_r5/ep8}, \nolinkurl{L/astra_r7/ep78b} & \nolinkurl{audits_r2_r6/astra_r6_auditStruct.report.md}; \nolinkurl{V/astra_r7v/v_ep78b}\\
Ex.~\ref{ex:cb}, \ref{ex:heawood}, Conj.~\ref{conj:cbhigh} & \nolinkurl{L/astra_r10/cb}; \nolinkurl{referee_data/r10/petersen_cb.json} & \nolinkurl{V/astra_r10v/v_cb}\\
App.~\ref{app:routes} & \nolinkurl{L/astra_r4/idea*}, \nolinkurl{L/astra_r5/ideas}, \nolinkurl{L/astra_r10/td} & \nolinkurl{verify_round_branch_r5.py}; \nolinkurl{V/astra_r10v/v_td}\\
Thm.~G / App.~\ref{app:window} (R35--R40) & \nolinkurl{ROUND_BRANCH_R35_CERTIFICATE.json}--\texttt{R40}; session reports under \nolinkurl{L/astra_r35}--\texttt{r40} & \nolinkurl{L/astra_r35v}--\texttt{r40v} verdicts; overall CWF\\
\bottomrule
\end{longtable}
\endgroup


\begin{thebibliography}{99}

\bibitem{cs}
E.~Camby and O.~Schaudt, \emph{A new characterization of $P_k$-free graphs}, Algorithmica \textbf{75} (2016), 205--217. \url{https://doi.org/10.1007/s00453-015-9989-6}; arXiv:1402.7213.

\bibitem{chen}
H.~Chen, \emph{Partial results on the WOWII-133 Graffiti conjecture}, Zenodo record 22054835, 2026 (version cited; concept DOI 10.5281/zenodo.22054834). \url{https://doi.org/10.5281/zenodo.22054835}.

\bibitem{wowii}
E.~DeLaVina, \emph{Written on the Wall II (Conjectures of Graffiti.pc)}, \url{http://cms.dt.uh.edu/faculty/delavinae/research/wowII/}, as referenced in~\cite{fc133}.

\bibitem{ess}
P.~Erd\H{o}s, M.~Saks and V.~T.~S\'os, \emph{Maximum induced trees in graphs}, J.~Combin. Theory Ser.~B \textbf{41} (1986), 61--79.

\bibitem{fc133}
Formal Conjectures (Google DeepMind), file \nolinkurl{FormalConjectures/WrittenOnTheWallII/GraphConjecture133.lean}, \url{https://github.com/google-deepmind/formal-conjectures}, accessed 17 October 2026.

\bibitem{fcpr}
Formal Conjectures, pull request \#5727, ``drop the floor on \texttt{l G}'', \url{https://github.com/google-deepmind/formal-conjectures/pull/5727}, merged 13 September 2026.

\bibitem{mathlib}
The mathlib Community, \emph{The Lean mathematical library}, in: Proceedings of the 9th ACM SIGPLAN International Conference on Certified Programs and Proofs (CPP 2020), 367--381.

\bibitem{lean4}
L.~de Moura and S.~Ullrich, \emph{The Lean 4 theorem prover and programming language}, in: Automated Deduction -- CADE 28, Lecture Notes in Comput. Sci. \textbf{12699}, Springer, 2021, 625--635.

\bibitem{nauty}
B.~D.~McKay and A.~Piperno, \emph{Practical graph isomorphism, II}, J.~Symbolic Comput. \textbf{60} (2014), 94--112.

\bibitem{seinsche}
D.~Seinsche, \emph{On a property of the class of $n$-colorable graphs}, J.~Combin. Theory Ser.~B \textbf{16} (1974), 191--193.

\end{thebibliography}

\end{document}
